% 第14章 函数项级数与幂级数——题目解答
% 覆盖：79 道主问题；含所有明确编号小问共 130 个编号项。
% 仅解答原章中的练习、思考题与参考题；原书已有解答的普通例题不重复。

\chapter*{第14章\quad 函数项级数与幂级数习题解答}
\addcontentsline{toc}{chapter}{第14章\quad 函数项级数与幂级数习题解答}
\subsection*{14.1.3 练习题}

\paragraph{14.1.3 练习题 1}
\begin{xsolution}
逐小题讨论极限函数以及误差的上确界。

\textbf{(1)} 对每个固定的 $x$，
\[
  S_n(x)=\frac{x}{1+n^2x^2}\to0.
\]
令 $y=n\abs{x}$，则
\[
  \sup_{x\in\RR}\abs{S_n(x)}
  =\frac1n\sup_{y\geqslant0}\frac{y}{1+y^2}
  =\frac1{2n}\to0.
\]
故在 $(-\infty,+\infty)$ 上一致收敛于 $0$。

\textbf{(2)} 对每个固定的 $x$，仍有 $S_n(x)\to0$。但
\[
  \sup_{x\in\RR}\abs{S_n(x)}
  =\sup_{y\geqslant0}\frac{y}{1+y^2}=\frac12,
  \qquad y=n\abs{x}.
\]
故不一致收敛。

\textbf{(3)}
\[
  S_n(x)=\frac{n+x^2}{nx}=\frac1x+\frac{x}{n},\qquad x>0,
\]
故逐点收敛于 $S(x)=1/x$。但
\[
  \sup_{x>0}\abs{S_n(x)-S(x)}
  =\sup_{x>0}\frac{x}{n}=+\infty,
\]
所以在 $(0,+\infty)$ 上不一致收敛。

\textbf{(4)} 对固定 $x>0$，令 $y=x/n$，则 $y\to0+$，且
\[
  S_n(x)=y\ln y\to0.
\]
在 $x\in(0,1)$ 时，充分大的 $n$ 满足 $0<y<1/n<\ee^{-1}$，函数 $-y\ln y$ 在该区间递增，因而
\[
  \sup_{0<x<1}\abs{S_n(x)}
  =\frac{\ln n}{n}\to0.
\]
故在 $(0,1)$ 上一致收敛。若 $x\in(0,+\infty)$，则对任意固定 $n$，$\abs{(x/n)\ln(x/n)}$ 随 $x\to+\infty$ 无界，因此不可能一致收敛于 $0$。

\textbf{(5)} 对固定 $x\geqslant0$，由 $\sin u/u\to1$ 得
\[
  n\sin\frac{x}{n}\to x.
\]
在 $[0,a]$ 上利用 $0\leqslant u-\sin u\leqslant u^3/6$，得到
\[
  \sup_{0\leqslant x\leqslant a}
  \abs{n\sin\frac{x}{n}-x}
  \leqslant\frac{a^3}{6n^2}\to0,
\]
故一致收敛于 $x$。在 $(0,+\infty)$ 上取 $x_n=n\pi$，则
\[
  \abs{n\sin(x_n/n)-x_n}=n\pi,
\]
故不一致收敛。

\textbf{(6)} 当 $x\geqslant0$ 时，极限函数为 $0$，且
\[
  0\leqslant \frac1n\ln(1+\ee^{-nx})\leqslant\frac{\ln2}{n},
\]
故在 $[0,+\infty)$ 上一致收敛于 $0$。当 $x<0$ 时，
\[
  \frac1n\ln(1+\ee^{-nx})
  =-x+\frac1n\ln(1+\ee^{nx}),
\]
而 $0<\ee^{nx}<1$，故
\[
  \sup_{x<0}\abs{S_n(x)+x}\leqslant\frac{\ln2}{n}\to0.
\]
所以在 $(-\infty,0)$ 上一致收敛于 $-x$。

\textbf{(7)} 当 $1/t\leqslant\abs{x}\leqslant t$ 时，
\[
  \abs{x^n+x^{-n}}\leqslant2t^n.
\]
于是通项绝对值不超过
\[
  M_n=\frac{2n^8t^n}{\sqrt{n!}}.
\]
并且
\[
  \frac{M_{n+1}}{M_n}
  =t\left(1+\frac1n\right)^8\frac1{\sqrt{n+1}}\to0.
\]
故 $\sum M_n$ 收敛，由 Weierstrass 判别法知原级数在给定集合上绝对一致收敛。

\textbf{(8)} 先考察
\[
  B_N(x)=\sum_{k=1}^N\sin x\sin kx.
\]
当 $x$ 不是 $2\pi$ 的整数倍时，
\[
  \sum_{k=1}^N\sin kx
  =\frac{\cos(x/2)-\cos((N+1/2)x)}{2\sin(x/2)},
\]
从而
\[
  \abs{B_N(x)}
  \leqslant2.
\]
在 $x=2m\pi$ 时 $B_N(x)=0$，故上述一致有界性对一切 $x>0$ 成立。另一方面
\[
  a_n(x)=\frac1{\sqrt{n+x}}
\]
关于 $n$ 单调减少，且
\[
  \sup_{x>0}a_n(x)\leqslant\frac1{\sqrt n}\to0.
\]
由一致 Dirichlet 判别法，原级数在 $(0,+\infty)$ 上一致收敛。

\textbf{(9)} 设 $0\leqslant x\leqslant b$。利用
\[
  \ln(1+u)\geqslant u-\frac{u^2}{2}\qquad(u\geqslant0),
\]
有
\[
  \left(1+\frac xn\right)^n
  =\exp\left(n\ln\left(1+\frac xn\right)\right)
  \geqslant \ee^{x-x^2/(2n)}.
\]
故
\[
  0\leqslant \ee^x-\left(1+\frac xn\right)^n
  \leqslant \ee^x\left(1-\ee^{-x^2/(2n)}\right)
  \leqslant \frac{\ee^b b^2}{2n}.
\]
因此级数通项满足
\[
  0\leqslant \frac1n\left[\ee^x-\left(1+\frac xn\right)^n\right]
  \leqslant\frac{\ee^b b^2}{2n^2},
\]
故在任意有限闭区间 $[0,b]$ 上一致收敛。对 $[0,+\infty)$，固定 $n$ 时
\[
  \frac1n\left[\ee^x-\left(1+\frac xn\right)^n\right]\to+\infty
  \qquad(x\to+\infty),
\]
所以通项甚至不一致趋于 $0$，级数在 $[0,+\infty)$ 上不一致收敛。
\end{xsolution}

\paragraph{14.1.3 练习题 2}
\begin{xsolution}
记
\[
  u_n(x)=\frac{x^n}{1+x^{2n}}.
\]
若 $\abs{x}<1$，则 $\abs{u_n(x)}\leqslant\abs{x}^n$；若 $\abs{x}>1$，则
\[
  \abs{u_n(x)}\leqslant\abs{x}^{-n}.
\]
故在 $\abs{x}\ne1$ 时绝对收敛。$x=1$ 时 $u_n=1/2$，$x=-1$ 时 $u_n=(-1)^n/2$，通项都不趋于 $0$，故发散。因此收敛域为
\[
  \RR\setminus\{-1,1\}.
\]

若 $0<r<1$，则在 $\abs{x}\leqslant r$ 上有 $\abs{u_n(x)}\leqslant r^n$；若 $R>1$，则在 $\abs{x}\geqslant R$ 上有 $\abs{u_n(x)}\leqslant R^{-n}$。因此在任何与 $\{\pm1\}$ 保持正距离的闭子集上均一致收敛，特别是在
\[
  [-r,r],\qquad (-\infty,-R]\cup[R,+\infty)
\]
上绝对一致收敛。

另一方面，在三个最大连通收敛区间 $(-\infty,-1)$、$(-1,1)$、$(1,+\infty)$ 上均不一致收敛，因为令 $x$ 从相应一侧趋于 $\pm1$ 时，$\abs{u_n(x)}$ 的上确界均为 $1/2$。所以不存在把某个 $\pm1$ 作为边界而又不与它保持正距离的最大一致收敛区间。
\end{xsolution}

\paragraph{14.1.3 练习题 3}
\begin{xsolution}
对每个固定 $x$，这是公比为 $-1/(1+x^2)$ 的几何级数；当 $x=0$ 时每项均为 $0$，当 $x\ne0$ 时绝对收敛。因此在整个实轴上处处绝对收敛。

原级数从第 $N+1$ 项起的余项为
\[
  R_N(x)=\frac{(-1)^N x^2}{x^2+2}\frac1{(1+x^2)^N}.
\]
令 $t=x^2\geqslant0$，则
\[
  \abs{R_N(x)}
  =\frac{t}{t+2}(1+t)^{-N}
  \leqslant\frac{t}{(1+t)^{N+1}}.
\]
右端在 $t=1/N$ 附近取得最大量级，且
\[
  \sup_{t\geqslant0}\frac{t}{(1+t)^{N+1}}
  =\frac1N\left(1+\frac1N\right)^{-N-1}\to0.
\]
故原级数在 $\RR$ 上一致收敛。

取绝对值后的级数在 $x\ne0$ 时的和为
\[
  x^2\sum_{n=1}^{\infty}\frac1{(1+x^2)^n}=1,
\]
而在 $x=0$ 时和为 $0$。其和函数在 $0$ 处不连续，而各项连续，所以绝对值级数不可能在 $\RR$ 上一致收敛。故原级数绝对收敛且一致收敛，但不绝对一致收敛。
\end{xsolution}

\paragraph{14.1.3 练习题 4}
\begin{xsolution}
令
\[
  q=\frac{a}{2\ln^2 2}<1.
\]
对 $\abs{x}<a$ 及 $n\geqslant2$，有
\[
  \abs{\frac{x}{n\ln^2n}}\leqslant q<1.
\]
对 $\abs{u}\leqslant q$，由中值定理
\[
  \abs{\ln(1+u)}\leqslant\frac{\abs{u}}{1-q}.
\]
于是
\[
  \abs{\ln\left(1+\frac{x}{n\ln^2n}\right)}
  \leqslant\frac{a}{1-q}\frac1{n\ln^2n}.
\]
而 $\sum_{n=2}^{\infty}1/(n\ln^2n)$ 收敛，故由 Weierstrass 判别法，所给函数项级数在 $(-a,a)$ 上一致收敛。
\end{xsolution}

\paragraph{14.1.3 练习题 5}
\begin{xsolution}
固定 $x>0$，由 $\abs{\sin u}\leqslant\abs{u}$，
\[
  \abs{2^n\sin\frac1{3^nx}}
  \leqslant\frac1x\left(\frac23\right)^n,
\]
故级数绝对收敛。

但取
\[
  x_n=\frac{2}{\pi3^n}>0,
\]
则第 $n$ 项为
\[
  2^n\sin\frac1{3^n x_n}=2^n.
\]
因此通项不一致趋于 $0$，原级数在 $(0,+\infty)$ 上不一致收敛。
\end{xsolution}

\paragraph{14.1.3 练习题 6}
\begin{xsolution}
对 $0\leqslant t\leqslant1$，有
\[
  0\leqslant\sin\pi t\leqslant\pi(1-t).
\]
故对 $x\in[0,1]$，
\[
  \abs{\int_0^x t^n\sin\pi t\,\dd t}
  \leqslant\pi\int_0^1 t^n(1-t)\,\dd t
  =\frac{\pi}{(n+1)(n+2)}.
\]
右端关于 $n$ 的级数收敛，所以由 Weierstrass 判别法，原函数项级数在 $[0,1]$ 上一致收敛。
\end{xsolution}

\paragraph{14.1.3 练习题 7}
\begin{xsolution}
由于 $\{f_n\}$ 在 $(a,b)$ 上一致收敛，它在该区间上一致满足 Cauchy 条件。给定 $\varepsilon>0$，存在 $N$，使对 $n,m\geqslant N$ 和所有 $x\in(a,b)$，有
\[
  \abs{f_n(x)-f_m(x)}<\varepsilon.
\]
令 $x\to a+$，利用 $f_n,f_m$ 在 $a$ 处连续，得
\[
  \abs{f_n(a)-f_m(a)}\leqslant\varepsilon.
\]
所以数列 $\{f_n(a)\}$ 为 Cauchy 数列，因而收敛。同理 $\{f_n(b)\}$ 也收敛。

再把上面的不等式分别取极限 $x\to a+$、$x\to b-$，可知对 $n,m\geqslant N$，在端点也有同样的 Cauchy 估计。因此
\[
  \sup_{x\in[a,b]}\abs{f_n(x)-f_m(x)}\leqslant\varepsilon,
\]
故 $\{f_n\}$ 在 $[a,b]$ 上一致收敛。
\end{xsolution}

\paragraph{14.1.3 练习题 8}
\begin{xsolution}
对任意固定 $x\in(a,b)$，当 $n$ 充分大时 $x\pm1/n\in(a,b)$，且由微分中值定理
\[
  F_n(x)
  =\frac{f(x+1/n)-f(x-1/n)}{2/n}\to f'(x).
\]
所以函数列逐点收敛于 $f'$。

现取任意内闭区间 $[\alpha,\beta]\subset(a,b)$。取 $\eta>0$ 使
\[
  [\alpha-\eta,\beta+\eta]\subset(a,b).
\]
当 $n>1/\eta$ 时，
\[
  F_n(x)=\frac n2\int_{x-1/n}^{x+1/n}f'(t)\,\dd t.
\]
因为 $f'$ 在闭区间 $[\alpha-\eta,\beta+\eta]$ 上一致连续，若记其连续模为 $\omega(\delta)$，则
\[
  \abs{F_n(x)-f'(x)}
  \leqslant\frac n2\int_{x-1/n}^{x+1/n}\abs{f'(t)-f'(x)}\,\dd t
  \leqslant\omega(1/n)\to0,
\]
且估计与 $x\in[\alpha,\beta]$ 无关。因此在每个内闭区间上一致收敛，即在 $(a,b)$ 上内闭一致收敛。

题面所给表达式在靠近端点处对较小的 $n$ 可能没有定义；上述结论只涉及充分大的尾列，这不影响点态极限与内闭一致收敛的结论。
\end{xsolution}

\paragraph{14.1.3 练习题 9}
\begin{xsolution}
不妨设所有 $u_n$ 都单调增加；单调减少时完全相同。对任意 $x\in[a,b]$，
\[
  \abs{u_n(x)}\leqslant\max\{\abs{u_n(a)},\abs{u_n(b)}\}
  \leqslant\abs{u_n(a)}+\abs{u_n(b)}.
\]
由题设，两端点处的两个绝对值级数都收敛，故
\[
  \sum_{n=1}^{\infty}\bigl(\abs{u_n(a)}+\abs{u_n(b)}\bigr)<\infty.
\]
由 Weierstrass 判别法，$\sum u_n(x)$ 在 $[a,b]$ 上绝对一致收敛，因而一致收敛。
\end{xsolution}

\paragraph{14.1.3 练习题 10}
\begin{xsolution}
由连续函数列的一致极限仍连续，$f\in C[a,b]$。题设 $f$ 在 $[a,b]$ 上无零点，所以
\[
  m=\min_{x\in[a,b]}\abs{f(x)}>0.
\]
由 $f_n\to f$ 一致，存在 $N$，使 $n\geqslant N$ 时
\[
  \sup_{[a,b]}\abs{f_n-f}<\frac m2.
\]
于是 $\abs{f_n(x)}\geqslant m/2$，从而对 $n\geqslant N$，
\[
  \abs{\frac1{f_n(x)}-\frac1{f(x)}}
  =\frac{\abs{f_n(x)-f(x)}}{\abs{f_n(x)f(x)}}
  \leqslant\frac2{m^2}\abs{f_n(x)-f(x)}.
\]
右端一致趋于 $0$，所以倒数函数列的尾列一致收敛于 $1/f$。题面写作“函数列 $\{1/f_n\}$”时已隐含各 $1/f_n$ 有定义；即使只由已给条件出发，也至少能保证充分大的 $n$ 时 $f_n$ 无零点，而有限项不影响一致收敛问题。
\end{xsolution}

\paragraph{14.1.3 练习题 11}
\begin{xsolution}
由一致收敛知函数列一致满足 Cauchy 条件。取 $N$ 使
\[
  \sup_{x\in I}\abs{f_n(x)-f_N(x)}<1,
  \qquad n\geqslant N.
\]
因 $f_N$ 有界，设 $\sup_I\abs{f_N}\leqslant M_N$，则对 $n\geqslant N$，
\[
  \abs{f_n(x)}\leqslant M_N+1.
\]
而前 $N-1$ 个函数各自有界。令
\[
  M=\max\{M_1,\ldots,M_{N-1},M_N+1\},
\]
即可得到 $\abs{f_n(x)}\leqslant M$ 对一切 $n$ 和 $x\in I$ 成立。因此函数列在 $I$ 上一致有界。
\end{xsolution}

\paragraph{14.1.3 练习题 12}
\begin{xsolution}
设 $f_n\to f$、$g_n\to g$ 均一致。由上一题，$\{f_n\}$ 与 $\{g_n\}$ 都一致有界；于是存在 $M>0$，使 $\abs{f_n(x)}\leqslant M$。同时 $g$ 也有界，设 $\abs{g(x)}\leqslant M$，必要时增大 $M$ 即可。于是
\[
\begin{aligned}
  \abs{f_ng_n-fg}
  &\leqslant\abs{f_n}\abs{g_n-g}+\abs{g}\abs{f_n-f}\\
  &\leqslant M\bigl(\abs{g_n-g}+\abs{f_n-f}\bigr).
\end{aligned}
\]
取上确界后右端趋于 $0$，故 $f_ng_n\to fg$ 一致。
\end{xsolution}

\paragraph{14.1.3 练习题 13}
\begin{xsolution}
对固定 $x$，$S_n(x)$ 是函数 $t\mapsto f(x+t)$ 在 $[0,1]$ 上的右端点 Riemann 和，因此
\[
  S_n(x)\to S(x):=\int_0^1f(x+t)\,\dd t.
\]

取任意闭区间 $[A,B]$。函数 $f$ 在 $[A,B+1]$ 上一致连续。记其连续模为
\[
  \omega(\delta)=\sup\{\abs{f(u)-f(v)}:\ u,v\in[A,B+1],\ \abs{u-v}\leqslant\delta\}.
\]
则 $\omega(\delta)\to0$。把积分按长度 $1/n$ 分段，有
\[
\begin{aligned}
  \abs{S_n(x)-S(x)}
  &\leqslant\sum_{k=1}^n\int_{(k-1)/n}^{k/n}
    \abs{f(x+k/n)-f(x+t)}\,\dd t\\
  &\leqslant\omega(1/n),
\end{aligned}
\]
该估计对所有 $x\in[A,B]$ 同时成立。因此在任意有界闭区间上一致收敛，即在 $(-\infty,+\infty)$ 上内闭一致收敛。
\end{xsolution}

\paragraph{14.1.3 练习题 14}
\begin{xsolution}
因 $S_1\in R[a,b]$，存在 $M>0$ 使 $\abs{S_1(x)}\leqslant M$。归纳证明
\[
  \abs{S_n(x)}\leqslant M\frac{\abs{x-a}^{n-1}}{(n-1)!},
  \qquad n\geqslant1.
\]
$n=1$ 时成立。若对 $n$ 成立，则
\[
\begin{aligned}
  \abs{S_{n+1}(x)}
  &\leqslant\int_a^x\abs{S_n(t)}\,\abs{\dd t}\\
  &\leqslant M\int_0^{\abs{x-a}}\frac{u^{n-1}}{(n-1)!}\,\dd u
  =M\frac{\abs{x-a}^n}{n!}.
\end{aligned}
\]
令 $L=b-a$，则
\[
  \sup_{x\in[a,b]}\abs{S_n(x)}
  \leqslant M\frac{L^{n-1}}{(n-1)!}\to0.
\]
故 $S_n$ 在 $[a,b]$ 上一致收敛于 $0$。
\end{xsolution}

\subsection*{14.2.4 练习题}

\paragraph{14.2.4 练习题 1}
\begin{xsolution}
对 $x\geqslant0$，有
\[
  0<\frac1{2^n n^x}\leqslant\frac1{2^n},
\]
且 $\sum 2^{-n}$ 收敛，所以原级数在 $x\geqslant0$ 上一致收敛。因而可以交换 $x\to0+$ 与求和次序：
\[
  \lim_{x\to0+}\sum_{n=1}^{\infty}\frac1{2^n n^x}
  =\sum_{n=1}^{\infty}\frac1{2^n}=1.
\]
\end{xsolution}

\paragraph{14.2.4 练习题 2}
\begin{xsolution}
对 $x\ne0$，
\[
  \abs{\left(x+\frac1n\right)^n}
  =\abs{x}^n\abs{1+\frac1{nx}}^n
  \sim\abs{x}^n\ee^{1/x}.
\]
故当 $\abs{x}<1$ 时收敛，当 $\abs{x}>1$ 时发散。$x=0$ 时为 $\sum n^{-n}$，收敛；$x=1$ 时通项趋于 $\ee$，$x=-1$ 时通项绝对值趋于 $\ee^{-1}$，均发散。因此定义域为
\[
  (-1,1).
\]

任取 $0<r<1$，再取 $q$ 使 $r<q<1$。充分大的 $n$ 满足对 $\abs{x}\leqslant r$ 有
\[
  \abs{x+1/n}\leqslant q,
\]
故级数尾项由几何级数 $\sum q^n$ 控制，于 $[-r,r]$ 一致收敛。因此 $S$ 在 $(-1,1)$ 上连续。

逐项求导得到
\[
  \frac{\dd}{\dd x}\left(x+\frac1n\right)^n
  =n\left(x+\frac1n\right)^{n-1}.
\]
在同一闭区间 $[-r,r]$ 上，充分大的 $n$ 时其绝对值不超过 $nq^{n-1}$，而 $\sum nq^{n-1}$ 收敛。因此导数级数在任意内闭区间上一致收敛，故
\[
  S'(x)=\sum_{n=1}^{\infty}n\left(x+\frac1n\right)^{n-1},
  \qquad -1<x<1.
\]
所以 $S$ 在定义域内连续且可微，事实上同理可知它在 $(-1,1)$ 内无限次可微。
\end{xsolution}

\paragraph{14.2.4 练习题 3}
\begin{xsolution}
把通项分为
\[
  \frac{x+n(-1)^n}{x^2+n^2}
  =\frac{x}{x^2+n^2}+(-1)^n\frac{n}{x^2+n^2}.
\]
固定 $M>0$。在 $\abs{x}\leqslant M$ 上，第一部分满足
\[
  \abs{\frac{x}{x^2+n^2}}\leqslant\frac{M}{n^2},
\]
故绝对一致收敛。对第二部分，当 $n>M$ 时
\[
  a_n(x)=\frac{n}{x^2+n^2}
\]
关于 $n$ 单调减少，且 $\sup_{\abs{x}\leqslant M}a_n(x)\leqslant1/n\to0$。由一致的交错级数判别法，第二部分也一致收敛。因此原级数在任意有界闭区间上一致收敛。

通项导数为
\[
  \frac{n^2-x^2-2nx(-1)^n}{(x^2+n^2)^2}.
\]
当 $\abs{x}\leqslant M$ 且 $n>M$ 时，
\[
  \abs{\frac{n^2-x^2-2nx(-1)^n}{(x^2+n^2)^2}}
  \leqslant\frac{1}{n^2}+\frac{M^2}{n^4}+\frac{2M}{n^3}.
\]
右端可求和，故导数级数在任意有界闭区间上一致收敛。由逐项求导定理，和函数在 $\RR$ 上可导，并且
\[
  S'(x)=\sum_{n=1}^{\infty}
  \frac{n^2-x^2-2nx(-1)^n}{(x^2+n^2)^2}.
\]
\end{xsolution}

\paragraph{14.2.4 练习题 4}
\begin{xsolution}
设 $S_n\to S$ 在 $\RR$ 上一致。给定 $\varepsilon>0$，取 $N$ 使
\[
  \sup_{x\in\RR}\abs{S_N(x)-S(x)}<\frac\varepsilon3.
\]
由于 $S_N$ 一致连续，存在 $\delta>0$，使 $\abs{x-y}<\delta$ 时
\[
  \abs{S_N(x)-S_N(y)}<\frac\varepsilon3.
\]
于是
\[
\begin{aligned}
  \abs{S(x)-S(y)}
  &\leqslant\abs{S(x)-S_N(x)}+\abs{S_N(x)-S_N(y)}+\abs{S_N(y)-S(y)}\\
  &<\varepsilon.
\end{aligned}
\]
故 $S$ 在 $(-\infty,+\infty)$ 上一致连续。
\end{xsolution}

\paragraph{14.2.4 练习题 5}
\begin{xsolution}
记 $f(x)=\lim_{n\to\infty}f_n(x)$。由局部一致收敛知 $f$ 在 $[a,b)$ 上连续；在端点 $b$ 处即使可能不连续，也不影响 Riemann 可积性。又由 $\abs{f_n(x)}\leqslant g(x)$ 取极限得 $\abs{f(x)}\leqslant g(x)$。因此 $f$ 有界且除至多端点 $b$ 外连续，从而 $f\in R[a,b]$。

设 $M=\sup_{[a,b]}g<\infty$。给定 $\varepsilon>0$，先取 $\delta>0$ 使
\[
  2M\delta<\frac\varepsilon2.
\]
在闭区间 $[a,b-\delta]$ 上 $f_n\to f$ 一致，因此存在 $N$，使 $n\geqslant N$ 时
\[
  \sup_{x\in[a,b-\delta]}\abs{f_n(x)-f(x)}
  <\frac{\varepsilon}{2(b-a)}.
\]
于是
\[
\begin{aligned}
  \abs{\int_a^b(f_n-f)\,\dd x}
  &\leqslant\int_a^{b-\delta}\abs{f_n-f}\,\dd x
  +\int_{b-\delta}^b\abs{f_n-f}\,\dd x\\
  &\leqslant\frac\varepsilon2+2M\delta<\varepsilon.
\end{aligned}
\]
故
\[
  \lim_{n\to\infty}\int_a^b f_n(x)\,\dd x
  =\int_a^b f(x)\,\dd x.
\]
\end{xsolution}

\paragraph{14.2.4 练习题 6}
\begin{xsolution}
先证必要性。若 $f$ 在 $x_0$ 连续，给定 $\varepsilon>0$。由 $f_n(x_0)\to f(x_0)$，可取 $N$ 使
\[
  \abs{f_N(x_0)-f(x_0)}<\frac\varepsilon3.
\]
又由 $f_N$ 和 $f$ 在 $x_0$ 连续，可取 $\delta>0$，使 $x\in O_\delta(x_0)$ 时
\[
  \abs{f_N(x)-f_N(x_0)}<\frac\varepsilon3,
  \qquad
  \abs{f(x)-f(x_0)}<\frac\varepsilon3.
\]
故
\[
  \abs{f_N(x)-f(x)}<\varepsilon.
\]

再证充分性。给定 $\varepsilon>0$，由题设取 $\delta_1>0$ 和 $N$，使
\[
  \abs{f_N(x)-f(x)}<\frac\varepsilon3,
  \qquad x\in O_{\delta_1}(x_0).
\]
特别在 $x=x_0$ 时也有
\[
  \abs{f_N(x_0)-f(x_0)}<\frac\varepsilon3.
\]
又因 $f_N$ 在 $x_0$ 连续，可取 $\delta_2>0$，使 $\abs{x-x_0}<\delta_2$ 时
\[
  \abs{f_N(x)-f_N(x_0)}<\frac\varepsilon3.
\]
取 $\delta=\min(\delta_1,\delta_2)$，则
\[
  \abs{f(x)-f(x_0)}<\varepsilon.
\]
故 $f$ 在 $x_0$ 连续。
\end{xsolution}

\paragraph{14.2.4 练习题 7}
\begin{xsolution}
若 $f_n\to f$ 一致，且 $x_n\to x_0$，则
\[
  \abs{f_n(x_n)-f(x_0)}
  \leqslant\sup_{[a,b]}\abs{f_n-f}+\abs{f(x_n)-f(x_0)}\to0,
\]
其中使用了 $f\in C[a,b]$。故条件必要。

反之，假设不一致收敛，则存在 $\varepsilon_0>0$、严格递增的指标 $n_k$ 和点 $y_k\in[a,b]$，使
\[
  \abs{f_{n_k}(y_k)-f(y_k)}\geqslant\varepsilon_0.
\]
由 $[a,b]$ 紧致，可取子列仍记为 $y_k$，使 $y_k\to x_0\in[a,b]$。构造数列
\[
  x_n=
  \begin{cases}
    y_k,&n=n_k,\\
    x_0,&n\ne n_k.
  \end{cases}
\]
则 $x_n\to x_0$。由题设应有 $f_n(x_n)\to f(x_0)$。特别
\[
  f_{n_k}(y_k)\to f(x_0).
\]
另一方面 $f(y_k)\to f(x_0)$，故
\[
  \abs{f_{n_k}(y_k)-f(y_k)}\to0,
\]
与其始终不小于 $\varepsilon_0$ 矛盾。因此必一致收敛。
\end{xsolution}

\paragraph{14.2.4 练习题 8}
\begin{xsolution}
对每个 $n$ 取 $f_n$ 的一个原函数 $F_n$，并用加常数的方法令
\[
  F_n(a)=0.
\]
由 $f_n\to f$ 一致，导函数列 $\{f_n\}$ 一致满足 Cauchy 条件。对任意 $m,n$，由微分中值定理，
\[
  \abs{F_n(x)-F_m(x)}
  \leqslant(b-a)\sup_{t\in[a,b]}\abs{f_n(t)-f_m(t)}.
\]
所以 $\{F_n\}$ 在 $[a,b]$ 上一致 Cauchy，存在一致极限 $F$。

下面证明 $F'=f$。给定 $\varepsilon>0$，取一个固定的 $N$，使
\[
  \sup_{[a,b]}\abs{f-f_N}<\varepsilon.
\]
对任意 $x,y\in[a,b]$，由中值定理及 $F_m\to F$，
\[
\begin{aligned}
  \abs{[F(y)-F_N(y)]-[F(x)-F_N(x)]}
  &=\lim_{m\to\infty}\abs{[F_m(y)-F_N(y)]-[F_m(x)-F_N(x)]}\\
  &\leqslant\varepsilon\abs{y-x}.
\end{aligned}
\]
因此对 $y\ne x$，
\[
  \abs{\frac{F(y)-F(x)}{y-x}-\frac{F_N(y)-F_N(x)}{y-x}}
  \leqslant\varepsilon.
\]
令 $y\to x$，并利用 $F_N'(x)=f_N(x)$，再结合 $\abs{f_N(x)-f(x)}<\varepsilon$，可得
\[
  \limsup_{y\to x}\abs{\frac{F(y)-F(x)}{y-x}-f(x)}\leqslant2\varepsilon.
\]
因 $\varepsilon$ 任意，$F'(x)=f(x)$。故 $f$ 在 $[a,b]$ 上有原函数 $F$。
\end{xsolution}

\subsection*{14.3.2 思考题}

\paragraph{14.3.2 思考题 1}
\begin{xsolution}
设幂级数中心为 $x_0$。既然在 $a,b$ 两点都收敛，则它的收敛区间必包含 $[a,b]$。对任意 $x\in(a,b)$，总可找到端点 $c\in\{a,b\}$ 使
\[
  \abs{x-x_0}<\abs{c-x_0},
\]
而级数在 $c$ 收敛。由 Abel 第一定理，级数在 $x$ 绝对收敛。因此在 $(a,b)$ 上处处绝对收敛；端点 $a,b$ 处只由题设保证收敛，可能是条件收敛。

再证一致收敛。若 $x_0\in[a,b]$，分别在 $[a,x_0]$ 和 $[x_0,b]$ 应用 Abel 判别法，可得两边均一致收敛；若 $x_0$ 在 $[a,b]$ 外，则取离 $x_0$ 较远的那个端点，级数在从 $x_0$ 到该端点的整个线段上一致收敛，从而在其子区间 $[a,b]$ 上一致收敛。故该级数在 $[a,b]$ 上收敛且一致收敛，并在开区间 $(a,b)$ 上绝对收敛。
\end{xsolution}

\paragraph{14.3.2 思考题 2}
\begin{xsolution}
记
\[
  \rho_a=\limsup_{n\to\infty}\sqrt[n]{\abs{a_n}}=\frac1{R_1},
  \qquad
  \rho_b=\limsup_{n\to\infty}\sqrt[n]{\abs{b_n}}=\frac1{R_2}.
\]

对和系数，有
\[
  \limsup\sqrt[n]{\abs{a_n+b_n}}
  \leqslant\max(\rho_a,\rho_b),
\]
故
\[
  R_{a+b}\geqslant\min(R_1,R_2).
\]
若 $R_1\ne R_2$，则实际上必有
\[
  R_{a+b}=\min(R_1,R_2).
\]
例如若 $R_1<R_2$ 而 $R_{a+b}>R_1$，则
\[
  \sum a_nx^n=\sum(a_n+b_n)x^n-\sum b_nx^n
\]
将在某个半径大于 $R_1$ 的圆内收敛，矛盾。当 $R_1=R_2$ 时可能发生系数抵消，故半径可以严格增大，甚至变成 $+\infty$，如取 $b_n=-a_n$。

对逐项乘积系数，
\[
  \limsup\sqrt[n]{\abs{a_nb_n}}
  \leqslant\rho_a\rho_b.
\]
因此在 $0<R_1,R_2<+\infty$ 时
\[
  R_{ab}\geqslant R_1R_2.
\]
等号不一定成立。例如令 $a_n$ 只在偶数 $n$ 非零，$b_n$ 只在奇数 $n$ 非零，则 $a_nb_n\equiv0$，乘积幂级数的收敛半径为 $+\infty$。若 $0<R_1,R_2<+\infty$，并且两个根值极限都存在，即 $\sqrt[n]{\abs{a_n}}\to1/R_1$、$\sqrt[n]{\abs{b_n}}\to1/R_2$，则 $\sqrt[n]{\abs{a_nb_n}}\to1/(R_1R_2)$，从而必有 $R_{ab}=R_1R_2$。当某个半径为 $0$ 或 $+\infty$ 时，应以 $\rho=1/R$ 的上极限形式理解上述结论，避免 $0\cdot\infty$ 的歧义。
\end{xsolution}

\paragraph{14.3.2 思考题 3}
\begin{xsolution}
设 $0<R<+\infty$。已知级数在整个开区间 $(x_0-R,x_0+R)$ 上一致收敛。由一致 Cauchy 条件，对任意 $\varepsilon>0$，存在 $N$，使 $m>n\geqslant N$ 时
\[
  \abs{\sum_{k=n}^{m}a_k(x-x_0)^k}<\varepsilon
\]
对所有 $\abs{x-x_0}<R$ 同时成立。固定 $m,n$，令 $x\to x_0+R$，由有限和的连续性得
\[
  \abs{\sum_{k=n}^{m}a_kR^k}\leqslant\varepsilon.
\]
所以 $\sum a_kR^k$ 满足 Cauchy 准则，在右端点收敛。同理令 $x\to x_0-R$，得
\[
  \abs{\sum_{k=n}^{m}a_k(-R)^k}\leqslant\varepsilon,
\]
故左端点也收敛。因此两个端点 $x=x_0\pm R$ 均收敛。
\end{xsolution}

\paragraph{14.3.2 思考题 4}
\begin{xsolution}
不能仅靠 Weierstrass 判别法证明 Abel 第二定理关于收敛端点的完整结论。

事实上，若只考察严格位于收敛半径内部的闭区间，则 Weierstrass 判别法当然有效。设收敛半径为 $R>0$，任取 $0<r<\rho<R$。因为幂级数在 $x=x_0+\rho$ 处收敛，所以数列 $\{a_n\rho^n\}$ 有界，即存在 $M>0$ 使
\[
  \abs{a_n}\rho^n\leqslant M.
\]
于是当 $\abs{x-x_0}\leqslant r$ 时，
\[
  \abs{a_n(x-x_0)^n}
  \leqslant M\left(\frac r\rho\right)^n,
\]
右边为收敛的几何级数。因此 M-判别法可以证明幂级数在每个严格内闭区间上绝对一致收敛。

但若某个端点属于收敛域而级数在那里只是条件收敛，则 M-判别法一般失效。例如
\[
  \sum_{n=1}^{\infty}\frac{(-1)^{n-1}}n x^n
\]
在 $x=1$ 处收敛，并且由 Abel 判别法可知它在 $[0,1]$ 上一致收敛；可是
\[
  \sum_{n=1}^{\infty}\sup_{0\leqslant x\leqslant1}
  \abs{\frac{(-1)^{n-1}}n x^n}
  =\sum_{n=1}^{\infty}\frac1n
\]
发散，因而不存在由这些逐项上确界构成的 Weierstrass 优级数。故 M-判别法只能直接处理收敛半径内部的绝对一致收敛，不能替代 Abel 定理在条件收敛端点处的论证。
\end{xsolution}

\paragraph{14.3.2 思考题 5}
\begin{xsolution}
\textbf{(1)} 一般不成立。取 $R=1$，并令
\[
  a_n=(n+1)(-1)^n,\qquad n=0,1,2,\ldots.
\]
相应幂级数的收敛半径为 $1$，且当 $\abs{x}<1$ 时
\[
  \sum_{n=0}^{\infty}a_nx^n
  =\sum_{n=0}^{\infty}(n+1)(-x)^n
  =\frac1{(1+x)^2}.
\]
因此
\[
  \int_0^1\left(\sum_{n=0}^{\infty}a_nx^n\right)\dd x
  =\int_0^1\frac{\dd x}{(1+x)^2}
  =\frac12,
\]
而题中右端变为
\[
  \sum_{n=0}^{\infty}\frac{a_n}{n+1}
  =\sum_{n=0}^{\infty}(-1)^n,
\]
它发散。故不能无条件把逐项积分公式直接延伸到收敛半径端点。

\textbf{(2)} 若
\[
  \sum_{n=0}^{\infty}\frac{a_n}{n+1}R^{n+1}
\]
收敛，则等式成立。令
\[
  F(r)=\sum_{n=0}^{\infty}\frac{a_n}{n+1}r^{n+1},
  \qquad 0\leqslant r<R.
\]
在 $r<R$ 时可以逐项积分，故
\[
  F(r)=\int_0^r\left(\sum_{n=0}^{\infty}a_nx^n\right)\dd x.
\]
又因为作为幂级数的 $F$ 在 $r=R$ 处收敛，由 Abel 第二定理，
\[
  \lim_{r\to R-}F(r)
  =\sum_{n=0}^{\infty}\frac{a_n}{n+1}R^{n+1}.
\]
所以
\[
  \int_0^R\left(\sum_{n=0}^{\infty}a_nx^n\right)\dd x
  =\sum_{n=0}^{\infty}\frac{a_n}{n+1}R^{n+1}.
\]
\end{xsolution}

\paragraph{14.3.2 思考题 6}
\begin{xsolution}
令 $y=x^2$。级数
\[
  \sum_{n=0}^{\infty}a_nx^{2n}
  =\sum_{n=0}^{\infty}a_ny^n
\]
收敛的条件为 $\abs{y}<R$，即 $\abs{x}<\sqrt R$。故作为关于 $x$ 的幂级数，其收敛半径为
\[
  \sqrt R,
\]
其中约定 $\sqrt0=0$、$\sqrt{+\infty}=+\infty$。
\end{xsolution}

\paragraph{14.3.2 思考题 7}
\begin{xsolution}
不能由 $\sum a_n$ 发散推出 $\lim_{x\to1-}S(x)$ 不存在。反例取
\[
  a_n=(-1)^n.
\]
则 $\sum(-1)^n$ 发散，但在 $\abs{x}<1$ 时
\[
  S(x)=\sum_{n=0}^{\infty}(-1)^nx^n=\frac1{1+x},
\]
所以
\[
  \lim_{x\to1-}S(x)=\frac12.
\]

若 $a_n\geqslant0$ 且 $\sum a_n$ 发散，则
\[
  \lim_{x\to1-}S(x)=+\infty.
\]
证明如下：给定任意 $M>0$，可取 $N$ 使 $\sum_{n=0}^Na_n>2M$。当 $x$ 充分接近 $1$ 时，$x^n>1/2$ 对 $0\leqslant n\leqslant N$ 同时成立，于是
\[
  S(x)\geqslant\sum_{n=0}^Na_nx^n>M.
\]
故极限为 $+\infty$，特别不存在有限极限。
\end{xsolution}

\subsection*{14.3.4 练习题}

\paragraph{14.3.4 练习题 1}
\begin{xsolution}
令
\[
  y=\frac{1-x}{1+x},
  \qquad
  A_n=\frac{1^n+2^n+\cdots+k^n}{n^2}.
\]
因为
\[
  k^n\leqslant1^n+2^n+\cdots+k^n\leqslant k\,k^n,
\]
故
\[
  \lim_{n\to\infty}\sqrt[n]{A_n}=k.
\]
所以关于 $y$ 的幂级数收敛半径为 $1/k$。当 $y=\pm1/k$ 时，
\[
  \abs{A_ny^n}
  \leqslant\frac{k}{n^2},
\]
故两个端点都绝对收敛。因此条件为
\[
  \abs{\frac{1-x}{1+x}}\leqslant\frac1k.
\]
因 $k>1$，平方化简得
\[
  (k^2-1)x^2-2(k^2+1)x+(k^2-1)\leqslant0.
\]
两个根为
\[
  \frac{k-1}{k+1},\qquad \frac{k+1}{k-1}.
\]
故收敛域为
\[
  \boxed{\left[\frac{k-1}{k+1},\frac{k+1}{k-1}\right]}.
\]
\end{xsolution}

\paragraph{14.3.4 练习题 2}
\begin{xsolution}
等差数列可写成
\[
  a_n=a_0+nd.
\]
若 $d=0$，则 $\sqrt[n]{\abs{a_n}}\to1$；若 $d\ne0$，仍有
\[
  \sqrt[n]{\abs{a_0+nd}}\to1.
\]
因为 $a_0\ne0$，级数不是恒为零的退化情形。由 Cauchy--Hadamard 公式，
\[
  R=\frac1{\limsup\sqrt[n]{\abs{a_n}}}=1.
\]
\end{xsolution}

\paragraph{14.3.4 练习题 3}
\begin{xsolution}
系数为
\[
  c_n=\frac{a^n}{1+a^{2n}}.
\]
若 $0<a<1$，则 $c_n\sim a^n$，故 $R=1/a$；若 $a=1$，则 $c_n=1/2$，故 $R=1$；若 $a>1$，则 $c_n\sim a^{-n}$，故 $R=a$。统一写成
\[
  \boxed{R=\max\left\{a,\frac1a\right\}}.
\]
\end{xsolution}

\paragraph{14.3.4 练习题 4}
\begin{xsolution}
记
\[
  c_n=\frac{(n!)^k}{(kn)!}.
\]
则
\[
  \frac{c_n}{c_{n+1}}
  =\frac{(kn+1)(kn+2)\cdots(kn+k)}{(n+1)^k}	o k^k.
\]
所以收敛半径为
\[
  R=k^k.
\]

还需讨论端点。由 Stirling 公式，
\[
\begin{aligned}
  c_n(k^k)^n
  &=\frac{(n!)^k k^{kn}}{(kn)!}\\
  &\sim\frac{(2\pi n)^{k/2}(n/\ee)^{kn}k^{kn}}
  {\sqrt{2\pi kn}(kn/\ee)^{kn}}
  =\frac{(2\pi n)^{k/2}}{\sqrt{2\pi kn}}.
\end{aligned}
\]
当 $k=1$ 时该量恒为 $1$；当 $k>1$ 时它甚至趋于 $+\infty$。因此在 $x=k^k$ 时通项不趋于 $0$；在 $x=-k^k$ 时通项绝对值相同，也不趋于 $0$。两个端点均发散，故收敛域为
\[
  \boxed{(-k^k,k^k)}.
\]
\end{xsolution}

\paragraph{14.3.4 练习题 5}
\begin{xsolution}
\textbf{(1)} 系数 $a_n=(1+1/n)^{n^2}$ 满足
\[
  \sqrt[n]{a_n}=\left(1+\frac1n\right)^n\to\ee,
\]
故 $R=\ee^{-1}$。在 $x=\pm\ee^{-1}$ 时，通项绝对值为
\[
  \left[\frac{(1+1/n)^n}{\ee}\right]^n.
\]
由
\[
  n\ln\left(1+\frac1n\right)=1-\frac1{2n}+\bigO\left(\frac1{n^2}\right)
\]
知上式趋于 $\ee^{-1/2}\ne0$。故两端点发散，收敛域为
\[
  \boxed{(-\ee^{-1},\ee^{-1})}.
\]

\textbf{(2)} 令
\[
  b_n=\frac1{n!}\left(\frac n\ee\right)^n.
\]
由 Stirling 公式 $b_n\sim1/\sqrt{2\pi n}$，故收敛半径为 $1$。在 $x=1$ 时原级数为 $\sum(-1)^n b_n$。又
\[
  \frac{b_{n+1}}{b_n}=\frac1\ee\left(1+\frac1n\right)^n<1,
\]
且 $b_n\to0$，故在 $x=1$ 条件收敛。在 $x=-1$ 时通项变成 $b_n>0$，与 $1/\sqrt n$ 同阶，故发散。因此收敛域为
\[
  \boxed{(-1,1]}.
\]

\textbf{(3)} 系数为
\[
  c_n=\frac{a^n+(-b)^n}{n},\qquad a,b>0.
\]
无论 $a=b$ 与否，
\[
  \limsup\sqrt[n]{\abs{c_n}}=\max(a,b),
\]
故 $R=1/\max(a,b)$。

若 $a>b$，在 $x=1/a$ 时通项含有 $1/n$ 的主部，发散；在 $x=-1/a$ 时主部为 $(-1)^n/n$，其余部分绝对收敛，所以条件收敛。故
\[
  \boxed{[-1/a,1/a)}.
\]
若 $b>a$，结论对称，得到
\[
  \boxed{(-1/b,1/b]}.
\]
若 $a=b$，系数在偶数项为 $2a^n/n$、奇数项为 $0$，两个端点都含调和子级数而发散，故
\[
  \boxed{(-1/a,1/a)}.
\]

\textbf{(4)}
\[
  a_n=\frac{3^{-\sqrt n}}{\sqrt{1+n^2}}
\]
满足 $\sqrt[n]{a_n}\to1$，故 $R=1$。在 $x=\pm1$ 时绝对值级数为
\[
  \sum_{n=0}^{\infty}\frac{3^{-\sqrt n}}{\sqrt{1+n^2}}.
\]
按 $m^2\leqslant n<(m+1)^2$ 分组，每组至多 $2m+1$ 项，且每项为 $\bigO(3^{-m}/m^2)$，所以第 $m$ 组为 $\bigO(3^{-m}/m)$，可求和。因此两端点均绝对收敛，收敛域为
\[
  \boxed{[-1,1]}.
\]

\textbf{(5)} 令 $H_n=1+1/2+\cdots+1/n$。有 $\sqrt[n]{H_n}\to1$，故 $R=1$。在 $x=\pm1$ 时通项绝对值为 $H_n\not\to0$，故两端点发散。因此收敛域为
\[
  \boxed{(-1,1)}.
\]

\textbf{(6)} 通项绝对值为
\[
  n^{n^2}\abs{x}^{n^3}
  =\left(n^{1/n}\abs{x}\right)^{n^3}.
\]
因为 $n^{1/n}\to1$，当 $\abs{x}<1$ 时充分大 $n$ 后有 $n^{1/n}\abs{x}<q<1$，故级数快速收敛；当 $\abs{x}\geqslant1$ 时通项不趋于 $0$。故收敛域为
\[
  \boxed{(-1,1)}.
\]

\textbf{(7)} 令 $y=x^2+x+1$。因为
\[
  \sin\frac1{3n}\sim\frac1{3n},
\]
关于 $y$ 的幂级数收敛半径为 $1$。在 $y=1$ 时与调和级数同阶而发散；在 $y=-1$ 时为交错型而收敛。但对实数 $x$，
\[
  y=\left(x+\frac12\right)^2+\frac34>0,
\]
所以只需 $y<1$，即
\[
  x(x+1)<0.
\]
因此收敛域为
\[
  \boxed{(-1,0)}.
\]

\textbf{(8)} 令
\[
  y=\frac{x}{3x+1},
  \qquad
  c_n=\frac{n-1}{n^2+1}.
\]
有 $c_n\sim1/n$，故关于 $y$ 的级数半径为 $1$。$y=1$ 时发散，$y=-1$ 时为交错级数而收敛。于是要求
\[
  \abs{\frac{x}{3x+1}}<1
\]
或取边界值 $-1$。前一不等式等价于
\[
  8x^2+6x+1>0,
\]
即 $x<-1/2$ 或 $x>-1/4$。其中 $x=-1/2$ 对应 $y=1$，应排除；$x=-1/4$ 对应 $y=-1$，应加入。因此收敛域为
\[
  \boxed{(-\infty,-1/2)\cup[-1/4,+\infty)}.
\]
\end{xsolution}

\paragraph{14.3.4 练习题 6}
\begin{xsolution}
\textbf{(1)} 级数半径为 $1$，且在 $x=\pm1$ 时由 $1/(n^2-1)$ 的可和性知均绝对收敛，所以收敛域为 $[-1,1]$。利用
\[
  \frac1{n^2-1}=\frac12\left(\frac1{n-1}-\frac1{n+1}\right),
\]
在 $-1<x\leqslant1$ 内计算
\[
\begin{aligned}
  S(x)
  &=\frac12x^2\sum_{k=1}^{\infty}\frac{(-1)^{k+1}x^k}{k}
   -\frac12\sum_{j=3}^{\infty}\frac{(-1)^{j-1}x^j}{j}\\
  &=\frac12\left[(x^2-1)\ln(1+x)+x-\frac{x^2}{2}\right].
\end{aligned}
\]
该表达式在 $x=-1$ 处取右极限 $-3/4$。因此
\[
  \boxed{S(x)=\frac12\left[(x^2-1)\ln(1+x)+x-\frac{x^2}{2}\right]}
\]
对 $-1<x\leqslant1$ 成立，并规定 $S(-1)=-3/4$。

\textbf{(2)} 半径为 $1$，且两个端点都绝对收敛，故收敛域为 $[-1,1]$。分解
\[
  \frac1{n(2n-1)}=-\frac1n+\frac2{2n-1},
\]
得
\[
\begin{aligned}
  S(x)
  &=-\sum_{n=1}^{\infty}\frac{(-1)^{n-1}x^{2n}}n
  +2x\sum_{n=1}^{\infty}\frac{(-1)^{n-1}x^{2n-1}}{2n-1}\\
  &=-\ln(1+x^2)+2x\arctan x.
\end{aligned}
\]
故
\[
  \boxed{S(x)=2x\arctan x-\ln(1+x^2),\quad -1\leqslant x\leqslant1}.
\]

\textbf{(3)} 令 $y=x^2$。级数为
\[
  2\sum_{n=1}^{\infty}(-1)^{n-1}y^n
  +\sum_{n=1}^{\infty}\frac{(-1)^{n-1}y^n}{n}.
\]
当 $\abs{x}<1$ 时
\[
  \boxed{S(x)=\frac{2x^2}{1+x^2}+\ln(1+x^2)}.
\]
在 $x=\pm1$ 时原通项不趋于 $0$，故收敛域为 $(-1,1)$。

\textbf{(4)} 利用
\[
  \sum_{n=1}^{\infty}n^2z^n=\frac{z(1+z)}{(1-z)^3},\qquad \abs z<1,
\]
取 $z=-x$，得
\[
  \boxed{S(x)=-\frac{x(1-x)}{(1+x)^3}},\qquad \abs{x}<1.
\]
在 $x=\pm1$ 时通项 $n^2(\pm1)^n$ 不趋于 $0$，故收敛域为 $(-1,1)$。

\textbf{(5)} 当 $\abs{x}<1/\sqrt2$ 时
\[
\begin{aligned}
  S(x)
  &=4\sum_{n=0}^{\infty}(2x^2)^n-\sum_{n=0}^{\infty}(x^2)^n\\
  &=\frac4{1-2x^2}-\frac1{1-x^2}
  =\boxed{\frac{3-2x^2}{(1-2x^2)(1-x^2)}}.
\end{aligned}
\]
在 $x=\pm1/\sqrt2$ 时通项不趋于 $0$，故收敛域为 $(-1/\sqrt2,1/\sqrt2)$。

\textbf{(6)} 令 $y=x-1$。当 $\abs y<1$ 时
\[
  \sum_{n=1}^{\infty}ny^n=\frac{y}{(1-y)^2}.
\]
故
\[
  \boxed{S(x)=\frac{x-1}{(2-x)^2}},
  \qquad 0<x<2.
\]
两个端点处通项不趋于 $0$，所以收敛域恰为 $(0,2)$。
\end{xsolution}

\paragraph{14.3.4 练习题 7}
\begin{xsolution}
级数的收敛半径为 $+\infty$，所以可以在整个实轴上逐项求导四次。得到
\[
\begin{aligned}
  u^{(4)}(x)
  &=\sum_{n=1}^{\infty}\frac{(4n)(4n-1)(4n-2)(4n-3)}{(4n)!}x^{4n-4}\\
  &=\sum_{n=1}^{\infty}\frac{x^{4n-4}}{(4n-4)!}
  =\sum_{m=0}^{\infty}\frac{x^{4m}}{(4m)!}=u(x).
\end{aligned}
\]
故 $u^{(4)}=u$ 在 $(-\infty,+\infty)$ 上成立。
\end{xsolution}

\paragraph{14.3.4 练习题 8}
\begin{xsolution}
记
\[
  a_n=\frac{(-1)^n}{(n!)^2 2^{2n}},
  \qquad J_0(x)=\sum_{n=0}^{\infty}a_nx^{2n}.
\]
该幂级数收敛半径为 $+\infty$，故可逐项求导。于是
\[
  xJ_0''(x)+J_0'(x)
  =\sum_{n=1}^{\infty}(2n)^2a_nx^{2n-1}.
\]
另一方面
\[
  xJ_0(x)=\sum_{n=1}^{\infty}a_{n-1}x^{2n-1}.
\]
而
\[
  a_{n-1}=-4n^2a_n=-(2n)^2a_n.
\]
故每个 $x^{2n-1}$ 的系数相消，得到
\[
  xJ_0''+J_0'+xJ_0=0.
\]
\end{xsolution}

\paragraph{14.3.4 练习题 9}
\begin{xsolution}
\textbf{(1)} 在 $[-1,1]$ 上有
\[
  \abs{\frac{x^n}{n^2\ln(1+n)}}
  \leqslant\frac1{n^2\ln(1+n)}.
\]
右端级数收敛，故原级数在 $[-1,1]$ 上绝对一致收敛。每项连续，所以
\[
  S\in C[-1,1].
\]

\textbf{(2)} 在 $(-1,1)$ 内可逐项求导：
\[
  S'(x)=\sum_{n=1}^{\infty}\frac{x^{n-1}}{n\ln(1+n)}.
\]
在 $x=-1$，右端成为
\[
  \sum_{n=1}^{\infty}\frac{(-1)^{n-1}}{n\ln(1+n)},
\]
其正项绝对值单调趋于 $0$，故交错收敛。由 Abel 第二定理用于导数幂级数，$S'(x)$ 当 $x\to-1+$ 时有有限极限。于是
\[
  S'_+(-1)=\sum_{n=1}^{\infty}\frac{(-1)^{n-1}}{n\ln(1+n)}
\]
有限。

在 $x=1$ 的左差商为
\[
\begin{aligned}
  \frac{S(1)-S(x)}{1-x}
  &=\sum_{n=1}^{\infty}
    \frac{1-x^n}{(1-x)n^2\ln(1+n)}\\
  &=\sum_{n=1}^{\infty}
    \frac{1+x+\cdots+x^{n-1}}{n^2\ln(1+n)}.
\end{aligned}
\]
当 $x\to1-$ 时，各项非负并递增到
\[
  \frac1{n\ln(1+n)}.
\]
而 $\sum1/[n\ln(1+n)]$ 发散，所以由单调性，对任意 $M$ 均可先取有限个项使极限和超过 $M$，再令 $x$ 足够接近 $1$。因此
\[
  S'_-(1)=+\infty.
\]
\end{xsolution}

\paragraph{14.3.4 练习题 10}
\begin{xsolution}
令 $y=2x$，则
\[
  f(x)=\sum_{n=1}^{\infty}\frac{y^n}{n^2}.
\]
在 $\abs y\leqslant1$ 上由 $\sum1/n^2$ 控制，故在 $[-1/2,1/2]$ 上绝对一致收敛，从而连续。

在 $\abs{x}<1/2$ 时可逐项求导：
\[
  f'(x)=\sum_{n=1}^{\infty}\frac{2^n}{n}x^{n-1}.
\]
在 $x=-1/2$ 处导数级数为
\[
  2\sum_{n=1}^{\infty}\frac{(-1)^{n-1}}n=2\ln2,
\]
由 Abel 定理可知右导数存在且等于 $2\ln2$。故 $f$ 在 $[-1/2,1/2)$ 上可导。

在 $x=1/2$ 处，左差商为
\[
  \frac{f(1/2)-f(x)}{1/2-x}
  =2\sum_{n=1}^{\infty}\frac{1-(2x)^n}{(1-2x)n^2}.
\]
当 $x\to1/2-$ 时第 $n$ 项增至 $2/n$，而调和级数发散。因此
\[
  f'_-(1/2)=+\infty,
\]
故在 $x=1/2$ 处不存在有限导数。
\end{xsolution}

\paragraph{14.3.4 练习题 11}
\begin{xsolution}
在 $0<x<1$，逐项求导得
\[
  f'(x)=\sum_{n=1}^{\infty}\frac{x^{n-1}}n
  =-\frac{\ln(1-x)}x.
\]
令
\[
  F(x)=f(x)+f(1-x)+\ln x\ln(1-x).
\]
则
\[
\begin{aligned}
  F'(x)
  &=-\frac{\ln(1-x)}x
  +\frac{\ln x}{1-x}
  +\frac{\ln(1-x)}x-\frac{\ln x}{1-x}=0.
\end{aligned}
\]
故 $F$ 在 $(0,1)$ 上为常数。令 $x\to0+$，利用
\[
  f(0)=0,\qquad f(1)=\sum_{n=1}^{\infty}\frac1{n^2}=\frac{\pi^2}{6},
  \qquad \ln x\ln(1-x)\to0,
\]
得到该常数为 $\pi^2/6$。因此
\[
  \boxed{f(x)+f(1-x)+\ln x\ln(1-x)=\frac{\pi^2}{6}},
  \qquad0<x<1.
\]
\end{xsolution}

\paragraph{14.3.4 练习题 12}
\begin{xsolution}
任取有界闭区间 $K=[-M,M]$。由于原幂级数收敛半径为 $+\infty$，部分和 $S_n\to S$ 在 $K$ 上一致。于是存在 $L>0$ 和 $N$，使对 $n\geqslant N$ 及 $x\in K$，
\[
  \abs{S_n(x)}\leqslant L,
  \qquad \abs{S(x)}\leqslant L.
\]
函数 $S$ 在整个实轴连续，故在 $[-L,L]$ 上一致连续。于是给定 $\varepsilon>0$，可取 $\delta>0$，使 $u,v\in[-L,L]$ 且 $\abs{u-v}<\delta$ 时
\[
  \abs{S(u)-S(v)}<\varepsilon.
\]
再由 $S_n\to S$ 在 $K$ 上一致，充分大 $n$ 时
\[
  \sup_{x\in K}\abs{S_n(x)-S(x)}<\delta.
\]
因此
\[
  \sup_{x\in K}\abs{S(S_n(x))-S(S(x))}<\varepsilon.
\]
所以 $S\circ S_n\to S\circ S$ 在任意有界闭区间上一致，即在 $\RR$ 上内闭一致收敛。
\end{xsolution}

\paragraph{14.3.4 练习题 13}
\begin{xsolution}
记
\[
  A(x)=\sum_{n=0}^{\infty}a_nx^n,
  \qquad
  B(x)=\sum_{n=0}^{\infty}S_nx^n,
  \qquad S_n=\sum_{k=0}^na_k.
\]

\textbf{(1)} 由 $a_n=S_n-S_{n-1}$（约定 $S_{-1}=0$），形式上有
\[
  A(x)=(1-x)B(x).
\]
若 $B$ 在 $\abs{x}<r$ 收敛，则 $A=(1-x)B$ 也在那里收敛，所以 $R_A\geqslant R_B$。

反过来，因为 $\sum a_n$ 在 $x=1$ 发散，所以 $R_A\leqslant1$；否则若 $R_A>1$，则在 $x=1$ 绝对收敛，矛盾。对任意 $\abs{x}<R_A\leqslant1$，
\[
  \frac{A(x)}{1-x}=A(x)\sum_{j=0}^{\infty}x^j
\]
为两个绝对收敛级数的 Cauchy 乘积，其 $x^n$ 系数正是 $S_n$。因此 $B$ 在 $\abs{x}<R_A$ 收敛，故 $R_B\geqslant R_A$。综上
\[
  \boxed{R_A=R_B}.
\]

\textbf{(2)} 现设 $a_n\geqslant0$ 且 $a_n/S_n\to0$。由 $\sum a_n$ 发散，$S_n\to+\infty$。并且
\[
  \frac{S_n}{S_{n-1}}
  =\frac1{1-a_n/S_n}\to1.
\]
因此
\[
  \sqrt[n]{S_n}\to1.
\]
例如取对数后用 Cesàro 平均即可得到这一结论。又 $0\leqslant a_n\leqslant S_n$，所以
\[
  \limsup\sqrt[n]{a_n}\leqslant1,
\]
从而 $R_A\geqslant1$。另一方面，若 $R_A>1$，则 $\sum a_n$ 在 $x=1$ 收敛，与题设矛盾。因此
\[
  \boxed{R_A=1}.
\]
\end{xsolution}

\subsection*{14.4.4 练习题}

\paragraph{14.4.4 练习题 1}
\begin{xsolution}
分别作几何级数展开。

按 $x$ 的乘幂展开：
\[
  \frac1{a-x}=\frac1a\frac1{1-x/a}
  =\sum_{n=0}^{\infty}\frac{x^n}{a^{n+1}},
  \qquad \abs{x}<\abs a.
\]

按 $x-b$ 的乘幂展开。若 $a\ne b$，则
\[
  \frac1{a-x}
  =\frac1{a-b}\frac1{1-(x-b)/(a-b)}
  =\sum_{n=0}^{\infty}\frac{(x-b)^n}{(a-b)^{n+1}},
\]
成立范围为
\[
  \abs{x-b}<\abs{a-b}.
\]
若 $a=b$，函数在展开点有极点，不存在通常意义下关于 $x-b$ 的幂级数展开式。

按 $1/x$ 的乘幂展开：
\[
  \frac1{a-x}
  =-\frac1x\frac1{1-a/x}
  =-\sum_{n=0}^{\infty}\frac{a^n}{x^{n+1}},
  \qquad \abs{x}>\abs a.
\]
\end{xsolution}

\paragraph{14.4.4 练习题 2}
\begin{xsolution}
令
\[
  z=\frac{x-1}{x+1}.
\]
则
\[
  x=\frac{1+z}{1-z}.
\]
对于 $x>0$ 有 $\abs z<1$，于是
\[
\begin{aligned}
  \ln x
  &=\ln(1+z)-\ln(1-z)\\
  &=2\sum_{n=0}^{\infty}\frac{z^{2n+1}}{2n+1}.
\end{aligned}
\]
所以
\[
  \boxed{
  \ln x=2\sum_{n=0}^{\infty}\frac1{2n+1}
  \left(\frac{x-1}{x+1}\right)^{2n+1}},
  \qquad x>0.
\]
\end{xsolution}

\paragraph{14.4.4 练习题 3}
\begin{xsolution}
由导函数序列在 $(-a,a)$ 上一致有界，存在 $M>0$，使
\[
  \abs{f^{(n)}(x)}\leqslant M,
  \qquad x\in(-a,a),\quad n=1,2,\ldots.
\]
固定 $x\in(-a,a)$。在 $0$ 与 $x$ 之间应用 Taylor 公式的 Lagrange 余项，
\[
  f(x)=\sum_{k=0}^{n}\frac{f^{(k)}(0)}{k!}x^k+r_n(x),
\]
且
\[
  \abs{r_n(x)}
  \leqslant M\frac{\abs{x}^{n+1}}{(n+1)!}\to0.
\]
因此对每个 $x\in(-a,a)$，
\[
  f(x)=\sum_{k=0}^{\infty}\frac{f^{(k)}(0)}{k!}x^k.
\]
故 $f$ 在 $(-a,a)$ 上可展开为其 Maclaurin 级数。
\end{xsolution}

\paragraph{14.4.4 练习题 4}
\begin{xsolution}
由导函数列在整个实轴上一致有界，存在 $M>0$，使
\[
  \abs{f^{(n)}(x)}\leqslant M,
  \qquad x\in\RR,\quad n=1,2,\ldots.
\]
因此与上一题相同，对每个固定 $x\in\RR$，Taylor 公式的余项满足
\[
  \abs{r_n(x)}
  \leqslant M\frac{\abs{x}^{n+1}}{(n+1)!}\to0,
\]
从而 $f$ 在整个实轴上等于其 Maclaurin 级数。

下面证明所有 Maclaurin 系数都为 $0$。从正数列 $x_n\to0$ 中可取一个严格递减子列，仍记为 $x_n$。由连续性，
\[
  f(0)=\lim_{n\to\infty}f(x_n)=0.
\]
对每个 $n$，由 Rolle 定理，在 $(x_{n+1},x_n)$ 中存在 $y_n$，使
\[
  f'(y_n)=0.
\]
显然 $y_n\to0$，故由 $f'$ 连续得到 $f'(0)=0$。再对 $\{y_n\}$ 重复同一论证，可得到趋于 $0$ 的 $f''$ 的零点列，从而 $f''(0)=0$。如此递推，
\[
  f^{(m)}(0)=0,\qquad m=0,1,2,\ldots.
\]
于是 $f$ 的 Maclaurin 级数恒为零，而该级数在整个实轴上等于 $f$，故
\[
  \boxed{f(x)\equiv0}.
\]
\end{xsolution}

\paragraph{14.4.4 练习题 5}
\begin{xsolution}
\textbf{(1)}
\[
  \sqrt{4-x}=2\left(1-\frac x4\right)^{1/2}
  =2\sum_{n=0}^{\infty}{1/2\choose n}\left(-\frac x4\right)^n.
\]
二项式级数在本题中于两端点也收敛，因此成立范围为
\[
  \boxed{-4\leqslant x\leqslant4}.
\]

\textbf{(2)} 对 $a>0$，
\[
  a^x=\ee^{x\ln a}
  =\boxed{\sum_{n=0}^{\infty}\frac{(\ln a)^n}{n!}x^n},
  \qquad x\in\RR.
\]

\textbf{(3)} 因
\[
  1+3x+2x^2=(1+x)(1+2x),
\]
所以在两展开式同时成立的范围内
\[
\begin{aligned}
  \ln(1+3x+2x^2)
  &=\ln(1+x)+\ln(1+2x)\\
  &=\sum_{n=1}^{\infty}\frac{(-1)^{n-1}(1+2^n)}n x^n.
\end{aligned}
\]
收敛半径由较近的奇点 $x=-1/2$ 决定，为 $1/2$。在 $x=1/2$ 时两对数级数均收敛，在 $x=-1/2$ 时函数本身无定义。因此成立范围为
\[
  \boxed{-\frac12<x\leqslant\frac12}.
\]

\textbf{(4)} 对 $\abs{x}<1$，利用正切倍角公式，
\[
  \arctan\frac{2x}{1-x^2}=2\arctan x,
\]
因为两边均在 $x=0$ 处取 $0$ 且在 $(-1,1)$ 上连续。故
\[
  \boxed{
  \arctan\frac{2x}{1-x^2}
  =2\sum_{n=0}^{\infty}\frac{(-1)^n}{2n+1}x^{2n+1}},
  \qquad -1<x<1.
\]
端点 $x=\pm1$ 处原函数中的分式无定义。

\textbf{(5)}
\[
  1+x+x^2+x^3+x^4=\frac{1-x^5}{1-x},
\]
所以
\[
  \frac1{1+x+x^2+x^3+x^4}
  =\frac{1-x}{1-x^5}
  =\boxed{\sum_{n=0}^{\infty}(x^{5n}-x^{5n+1})},
\]
成立于
\[
  \boxed{\abs{x}<1}.
\]
把它按通常的幂级数写成 $\sum_{m=0}^{\infty}a_mx^m$，则
\[
  a_{5n}=1,\qquad a_{5n+1}=-1,
\]
其余系数为 $0$。因此在 $x=1$ 或 $x=-1$ 时，通项序列都含有绝对值为 $1$ 的无穷多个项，不能趋于 $0$，故两个端点均不属于成立范围。

\textbf{(6)} 由 Cauchy 乘积，
\[
\begin{aligned}
  \frac{\ee^{x^2}}{1-x^2}
  &=\left(\sum_{k=0}^{\infty}\frac{x^{2k}}{k!}\right)
    \left(\sum_{j=0}^{\infty}x^{2j}\right)\\
  &=\boxed{\sum_{n=0}^{\infty}\left(\sum_{k=0}^{n}\frac1{k!}\right)x^{2n}},
\end{aligned}
\]
成立范围由 $1/(1-x^2)$ 的极点决定，为
\[
  \boxed{\abs{x}<1}.
\]

\textbf{(7)} 先有
\[
  \frac{\ee^x-1}{x}=\sum_{n=0}^{\infty}\frac{x^n}{(n+1)!},
\]
整个实轴上成立。逐项求导得
\[
  \boxed{
  \frac{\dd}{\dd x}\left(\frac{\ee^x-1}{x}\right)
  =\sum_{n=0}^{\infty}\frac{n+1}{(n+2)!}x^n},
  \qquad x\in\RR,
\]
其中左边在 $x=0$ 按可去间断的延拓理解。

\textbf{(8)}
\[
  \frac{1-\cos t}{t^2}
  =\sum_{n=1}^{\infty}\frac{(-1)^{n+1}}{(2n)!}t^{2n-2},
\]
该级数处处收敛。逐项积分得到
\[
  \boxed{
  \int_0^x\frac{1-\cos t}{t^2}\,\dd t
  =\sum_{n=1}^{\infty}
  \frac{(-1)^{n+1}}{(2n)!(2n-1)}x^{2n-1}},
  \qquad x\in\RR.
\]
\end{xsolution}

\paragraph{14.4.4 练习题 6}
\begin{xsolution}
\textbf{(1)} 令 $y=x-2$，则
\[
  \frac1{3+x}=\frac1{5+y}
  =\boxed{\frac15\sum_{n=0}^{\infty}\left(-\frac y5\right)^n},
  \qquad \abs{y}<5.
\]
即 $\abs{x-2}<5$。

\textbf{(2)} 令 $y=x+1$。有
\[
  3+2x+x^2=2+y^2,
\]
故
\[
\begin{aligned}
  \ln\frac1{3+2x+x^2}
  &=-\ln(2+y^2)\\
  &=-\ln2-\ln\left(1+\frac{y^2}{2}\right)\\
  &=\boxed{-\ln2+
  \sum_{n=1}^{\infty}\frac{(-1)^n}{n2^n}y^{2n}}.
\end{aligned}
\]
在 $y=\pm\sqrt2$ 时是交错调和型级数，也收敛，故成立范围为
\[
  \boxed{\abs{x+1}\leqslant\sqrt2}.
\]

\textbf{(3)} 令 $y=x-3$。则
\[
  \ln x=\ln3+\ln\left(1+\frac y3\right)
  =\boxed{\ln3+
  \sum_{n=1}^{\infty}\frac{(-1)^{n-1}}{n3^n}y^n}.
\]
在 $y=3$ 即 $x=6$ 时收敛，在 $y=-3$ 即 $x=0$ 时函数无定义且级数发散，故成立范围为
\[
  \boxed{0<x\leqslant6}.
\]

\textbf{(4)} 令 $y=x-1$。则
\[
  \frac{x-1}{x+1}=\frac y{2+y}
  =\frac y2\frac1{1+y/2}
  =\boxed{\sum_{n=1}^{\infty}(-1)^{n-1}\frac{y^n}{2^n}},
\]
成立于
\[
  \boxed{\abs{x-1}<2}.
\]

\textbf{(5)} 令 $y=x-\pi/6$。利用加法公式，
\[
  \sin x=\frac12\cos y+\frac{\sqrt3}{2}\sin y.
\]
所以
\[
  \boxed{
  \sin x=
  \frac12\sum_{n=0}^{\infty}\frac{(-1)^n y^{2n}}{(2n)!}
  +\frac{\sqrt3}{2}\sum_{n=0}^{\infty}\frac{(-1)^n y^{2n+1}}{(2n+1)!}},
\]
对所有 $y\in\RR$ 成立，即成立范围为整个实轴。
\end{xsolution}

\subsection*{14.5.2 第一组参考题}

\paragraph{第一组参考题 1}
\begin{xsolution}
取 Dirichlet 函数
\[
  D(x)=
  \begin{cases}
    1,&x\in\QQ,\\
    0,&x\notin\QQ,
  \end{cases}
  \qquad x\in[0,1],
\]
并令
\[
  S_n(x)=\frac1nD(x).
\]
由于有理数与无理数在任一区间中都稠密，每个 $S_n$ 在 $[0,1]$ 上处处不连续。但
\[
  \sup_{x\in[0,1]}\abs{S_n(x)}=\frac1n\to0,
\]
所以 $S_n$ 在 $[0,1]$ 上一致收敛于连续函数 $0$。
\end{xsolution}

\paragraph{第一组参考题 2}
\begin{xsolution}
对任意整数 $k\geqslant0$，形式逐项求导得到
\[
  \boxed{
  \zeta^{(k)}(x)=(-1)^k
  \sum_{n=1}^{\infty}\frac{(\ln n)^k}{n^x}},
  \qquad x>1.
\]
下面证明逐项求导合法。任取内闭区间 $[1+\delta,M]\subset(1,+\infty)$，其中 $\delta>0$。对 $x$ 在该区间中，
\[
  \frac{(\ln n)^k}{n^x}
  \leqslant\frac{(\ln n)^k}{n^{1+\delta}}.
\]
而级数
\[
  \sum_{n=2}^{\infty}\frac{(\ln n)^k}{n^{1+\delta}}
\]
由积分判别法收敛。因此任意阶导数所成级数都在每个内闭区间上一致收敛。逐次使用函数项级数逐项求导定理，即得上述公式，并说明 $\zeta$ 在 $(1,+\infty)$ 上无限次可微。
\end{xsolution}

\paragraph{第一组参考题 3}
\begin{xsolution}
设 Dirichlet 级数
\[
  \sum_{n=1}^{\infty}\frac{a_n}{n^x}
\]
在 $x=x_2$ 收敛。若 $x>x_2$，写成
\[
  \sum_{n=1}^{\infty}\left(\frac{a_n}{n^{x_2}}\right)
  \frac1{n^{x-x_2}}.
\]
第一因子所成级数收敛，第二因子关于 $n$ 单调趋于 $0$，故由 Dirichlet--Abel 判别法，该级数在每个 $x>x_2$ 也收敛。

同理，若级数在 $x=x_1$ 发散，则它不可能在任何 $x<x_1$ 收敛；否则由上一段会推出它在 $x_1$ 收敛。

因此收敛点集具有“向右封闭”的形状。定义
\[
  r=\inf\left\{x\in\RR:\sum_{n=1}^{\infty}\frac{a_n}{n^x}
  \text{ 收敛}\right\}.
\]
题设给出了一个发散点 $x_1$ 和一个收敛点 $x_2$，故 $r$ 为有限实数，并有
\[
  x<r\Longrightarrow\text{发散},
  \qquad
  x>r\Longrightarrow\text{收敛}.
\]
在边界点 $x=r$ 本题不作断言，既可能收敛也可能发散。
\end{xsolution}

\paragraph{第一组参考题 4}
\begin{xsolution}
若 $S_n(x)=x^nf(x)$ 在 $[0,1]$ 上一致收敛，则其极限函数连续。对 $0\leqslant x<1$，有 $x^nf(x)\to0$；而在 $x=1$ 处极限为 $f(1)$。连续性迫使 $f(1)=0$。

反过来设 $f(1)=0$。给定 $\varepsilon>0$。由 $f$ 在 $1$ 处连续，可取 $\delta\in(0,1)$，使 $1-\delta<x\leqslant1$ 时
\[
  \abs{f(x)}<\varepsilon.
\]
在 $[0,1-\delta]$ 上，设 $\abs{f(x)}\leqslant M$，则
\[
  \abs{x^nf(x)}\leqslant M(1-\delta)^n\to0
\]
一致成立。故充分大的 $n$ 时，在前一段误差小于 $\varepsilon$，后一段由 $\abs{x^nf(x)}\leqslant\abs{f(x)}<\varepsilon$。所以 $S_n\to0$ 在 $[0,1]$ 上一致收敛。

因此充分必要条件是
\[
  \boxed{f(1)=0}.
\]
\end{xsolution}

\paragraph{第一组参考题 5}
\begin{xsolution}
各项函数
\[
  f_n(x)=\frac{\abs{x-r_n}}{3^n}
\]
满足 $\abs{f_n(x)}\leqslant3^{-n}$，故级数在 $[0,1]$ 上一致收敛。每项连续，所以和函数 $f$ 处处连续。

若 $x_0$ 为无理数，则 $x_0\ne r_n$ 对所有 $n$。对 $h\ne0$，
\[
  \abs{
  \frac{\abs{x_0+h-r_n}-\abs{x_0-r_n}}h}
  \leqslant1.
\]
因此差商级数的第 $n$ 项绝对值不超过 $3^{-n}$。对每个固定 $n$，当 $h\to0$ 时差商趋于
\[
  \operatorname{sgn}(x_0-r_n).
\]
由一致控制可交换极限与求和，得到
\[
  f'(x_0)=\sum_{n=1}^{\infty}
  \frac{\operatorname{sgn}(x_0-r_n)}{3^n}.
\]
故每个无理点处可微。

若 $x_0=r_m$ 为有理点，把第 $m$ 项分离：
\[
  f(x)=\frac{\abs{x-r_m}}{3^m}+g(x).
\]
由上述同样的控制论证，$g$ 在 $r_m$ 可微。于是
\[
  f'_+(r_m)=g'(r_m)+3^{-m},
  \qquad
  f'_-(r_m)=g'(r_m)-3^{-m}.
\]
左右导数不同，因此在每个有理点处不可微。
\end{xsolution}

\paragraph{第一组参考题 6}
\begin{xsolution}
因为 $f$ 不恒为 $0$，存在 $x_0\in(0,1]$ 使 $f(x_0)\ne0$。由
\[
  f^{(j)}(0)=0,
  \qquad j=0,1,2,\ldots,
\]
对每个正整数 $n$ 在 $0$ 与 $x_0$ 之间使用 Taylor 公式的 Lagrange 余项，可得某个 $\xi_n\in(0,x_0)$，使
\[
  f(x_0)=\frac{f^{(n)}(\xi_n)}{n!}x_0^n.
\]
因此
\[
  \sup_{x\in[0,1]}\abs{f^{(n)}(x)}
  \geqslant\frac{n!\abs{f(x_0)}}{x_0^n}
  \geqslant n!\abs{f(x_0)}.
\]
由 $a_nf^{(n)}\to0$ 在 $[0,1]$ 上一致，
\[
  \abs{a_n}\sup_{[0,1]}\abs{f^{(n)}}\to0.
\]
故
\[
  n!\abs{a_n}\abs{f(x_0)}\to0,
\]
即
\[
  \boxed{n!a_n\to0}.
\]
\end{xsolution}

\paragraph{第一组参考题 7}
\begin{xsolution}
定义在 $\abs{x}<1$ 上的幂级数
\[
  A(x)=\sum_{n=1}^{\infty}a_nx^n,
  \qquad
  B(x)=\sum_{n=1}^{\infty}b_nx^n.
\]
由于原数项级数收敛，系数序列有界，所以这两个幂级数在 $\abs{x}<1$ 绝对收敛。其乘积可按 Cauchy 乘积计算：
\[
  A(x)B(x)=\sum_{n=1}^{\infty}c_nx^n=:C(x),
  \qquad \abs{x}<1.
\]
题设又知数项级数 $\sum c_n$ 收敛。由 Abel 第二定理，
\[
  \lim_{x\to1-}A(x)=A,
  \qquad
  \lim_{x\to1-}B(x)=B,
  \qquad
  \lim_{x\to1-}C(x)=C.
\]
在恒等式 $C(x)=A(x)B(x)$ 中令 $x\to1-$，得到
\[
  \boxed{C=AB}.
\]
\end{xsolution}

\paragraph{第一组参考题 8}
\begin{xsolution}
\textbf{(1)} 先记
\[
  F(x)=\frac12(\arcsin x)^2.
\]
由 $F'(x)=\arcsin x/\sqrt{1-x^2}$ 可直接验证
\[
  (1-x^2)F''(x)-xF'(x)=1.
\]
设 $F(x)=\sum_{n=1}^{\infty}c_nx^{2n}$。比较系数可得
\[
  c_1=\frac12,
  \qquad
  c_n=\frac{(2n-2)^2}{(2n)(2n-1)}c_{n-1},\quad n\geqslant2.
\]
归纳得到
\[
  c_n=\frac{4^{n-1}}{n^2{2n\choose n}}.
\]
因此
\[
  \frac12(\arcsin x)^2
  =\frac14\sum_{n=1}^{\infty}
  \frac{4^n}{n^2{2n\choose n}}x^{2n}.
\]
逐项求导，得到
\[
  \boxed{
  \frac{\arcsin x}{\sqrt{1-x^2}}
  =\sum_{n=1}^{\infty}
  \frac{2^{2n-1}}{n{2n\choose n}}x^{2n-1}},
  \qquad \abs{x}<1.
\]
这就是第 (1) 式。

\textbf{(2)} 再令
\[
  G(x)=\frac12\operatorname{arsinh}^2x
  =\frac12\ln^2(x+\sqrt{1+x^2}).
\]
因 $G'(x)=\operatorname{arsinh}x/\sqrt{1+x^2}$，可验证
\[
  (1+x^2)G''(x)+xG'(x)=1.
\]
相同的系数递推只多出交错符号，故
\[
  G(x)=\frac14\sum_{n=1}^{\infty}
  (-1)^{n-1}\frac{4^n}{n^2{2n\choose n}}x^{2n}.
\]
逐项求导得到第 (2) 式：
\[
  \boxed{
  \frac{\ln(x+\sqrt{1+x^2})}{\sqrt{1+x^2}}
  =\sum_{n=1}^{\infty}
  (-1)^{n-1}\frac{2^{2n-1}}{n{2n\choose n}}x^{2n-1}}.
\]
记
\[
  d_n=\frac{2^{2n-1}}{n{2n\choose n}}.
\]
则
\[
  \frac{d_{n+1}}{d_n}=\frac{2n}{2n+1}<1,
  \qquad d_n\sim\frac{\sqrt\pi}{2\sqrt n}\to0.
\]
所以在 $x=\pm1$ 时级数均为条件收敛的交错级数。由 Abel 第二定理，等式在两个端点仍成立，因此第 (2) 式的成立范围为
\[
  \boxed{-1\leqslant x\leqslant1}.
\]

\textbf{(3)} 第 (3) 式由 $2F(x)/x^2$ 直接得到：
\[
  \boxed{
  \left(\frac{\arcsin x}{x}\right)^2
  =\frac12\sum_{n=1}^{\infty}
  \frac{4^n}{n^2{2n\choose n}}x^{2n-2}}.
\]
利用 ${2n\choose n}\sim4^n/\sqrt{\pi n}$，端点 $x=\pm1$ 时系数绝对值为 $\bigO(n^{-3/2})$，故该展开式在 $[-1,1]$ 上成立；$x=0$ 取极限值 $1$。
\end{xsolution}

\paragraph{第一组参考题 9}
\begin{xsolution}
由于 $a_n\to L$，数列 $\{a_n\}$ 有界，所以
\[
  f(x)=\sum_{n=1}^{\infty}a_nx^n
\]
在 $\abs{x}<1$ 绝对收敛。

写 $a_n=L+b_n$，其中 $b_n\to0$。则
\[
  (1-x)f(x)=Lx+(1-x)\sum_{n=1}^{\infty}b_nx^n.
\]
只需证明后一项趋于 $0$。给定 $\varepsilon>0$，取 $N$ 使 $n>N$ 时 $\abs{b_n}<\varepsilon$。于是
\[
\begin{aligned}
  (1-x)\abs{\sum_{n=1}^{\infty}b_nx^n}
  &\leqslant(1-x)\sum_{n=1}^{N}\abs{b_n}x^n
  +(1-x)\sum_{n>N}\varepsilon x^n\\
  &\leqslant(1-x)C_N+\varepsilon x^{N+1}.
\end{aligned}
\]
令 $x\to1-$，右端上极限不超过 $\varepsilon$。因 $\varepsilon$ 任意，第二项趋于 $0$。故
\[
  \boxed{\lim_{x\to1-}(1-x)f(x)=L}.
\]
\end{xsolution}

\paragraph{第一组参考题 10}
\begin{xsolution}
由在整个实轴上一致收敛，$\{P_n\}$ 一致满足 Cauchy 条件。因此存在 $N$，使 $n,m\geqslant N$ 时
\[
  \sup_{x\in\RR}\abs{P_n(x)-P_m(x)}<1.
\]
于是多项式 $P_n-P_m$ 在整个实轴上有界。非恒定多项式在实轴上必无界，所以 $P_n-P_m$ 必为常数。

固定 $m=N$，可写
\[
  P_n(x)=P_N(x)+c_n,
  \qquad n\geqslant N.
\]
任取一个 $x_0$，由 $P_n(x_0)$ 收敛可知常数数列 $c_n$ 收敛于某个 $c$。因此
\[
  P(x)=P_N(x)+c,
\]
故 $P$ 是多项式。
\end{xsolution}

\paragraph{第一组参考题 11}
\begin{xsolution}
任选 $D+1$ 个互异点
\[
  0\leqslant x_0<x_1<\cdots<x_D\leqslant1.
\]
写
\[
  P_n(x)=\sum_{j=0}^{D}a_{n,j}x^j.
\]
由逐点收敛，向量
\[
  \bigl(P_n(x_0),\ldots,P_n(x_D)\bigr)
\]
收敛。它与系数向量 $\bigl(a_{n,0},\ldots,a_{n,D}\bigr)$ 之间由固定的 Vandermonde 矩阵联系。该矩阵可逆，因此每个系数数列 $a_{n,j}$ 都收敛于某个 $a_j$。

令
\[
  P(x)=\sum_{j=0}^{D}a_jx^j.
\]
则
\[
  \sup_{x\in[0,1]}\abs{P_n(x)-P(x)}
  \leqslant\sum_{j=0}^{D}\abs{a_{n,j}-a_j}\to0.
\]
所以 $P_n$ 在 $[0,1]$ 上一致收敛。
\end{xsolution}

\paragraph{第一组参考题 12}
\begin{xsolution}
设
\[
  f(x)=\sum_{n=0}^{\infty}a_nx^n
\]
在整个实轴上成立，其部分和记为 $P_N(x)$。若该幂级数在 $\RR$ 上一致收敛，则多项式序列 $\{P_N\}$ 在 $\RR$ 上一致收敛于 $f$。由上一题的同类结论（更直接地用第一组参考题 10），$f$ 必为多项式。

反之，若 $f$ 是多项式，则由 Maclaurin 展开的惟一性，只有有限个 $a_n$ 非零。此时部分和从某项起恒等于 $f$，当然在整个实轴上一致收敛。

故充分必要条件为
\[
  \boxed{f\text{ 是多项式}}.
\]
\end{xsolution}

\paragraph{第一组参考题 13}
\begin{xsolution}
由二项式系数
\[
  C(a)={a\choose2003}.
\]
因为 $2003$ 为奇数，
\[
\begin{aligned}
  C(-y-1)
  &={-y-1\choose2003}\\
  &=-\frac{(y+1)(y+2)\cdots(y+2003)}{2003!}.
\end{aligned}
\]
记
\[
  P(y)=\frac{(y+1)(y+2)\cdots(y+2003)}{2003!}.
\]
则
\[
  \frac{P'(y)}{P(y)}
  =\sum_{j=1}^{2003}\frac1{y+j}.
\]
所以被积函数正好为 $-P'(y)$。于是
\[
\begin{aligned}
  \int_0^1 C(-y-1)
  \left(\sum_{j=1}^{2003}\frac1{y+j}\right)\dd y
  &=-\int_0^1P'(y)\,\dd y\\
  &=P(0)-P(1).
\end{aligned}
\]
而
\[
  P(0)=1,
  \qquad
  P(1)=\frac{2\cdot3\cdots2004}{2003!}=2004.
\]
故积分值为
\[
  \boxed{-2003}.
\]
\end{xsolution}

\paragraph{第一组参考题 14}
\begin{xsolution}
记
\[
  T_{m,n}=\frac{m^2n}{3^m(n3^m+m3^n)}.
\]
交换 $m,n$ 后
\[
  T_{n,m}=\frac{n^2m}{3^n(m3^n+n3^m)}.
\]
两式分母中的括号相同，因此
\[
  T_{m,n}+T_{n,m}
  =\frac{mn}{3^{m+n}}.
\]
所有项均为正，可以自由交换与配对双重求和。设所求和为 $S$，则
\[
\begin{aligned}
  2S
  &=\sum_{m=1}^{\infty}\sum_{n=1}^{\infty}
    \frac{mn}{3^{m+n}}\\
  &=\left(\sum_{m=1}^{\infty}\frac m{3^m}\right)^2.
\end{aligned}
\]
而
\[
  \sum_{m=1}^{\infty}mr^m=\frac r{(1-r)^2},
  \qquad r=\frac13,
\]
故该一维级数之和为 $3/4$。所以
\[
  \boxed{S=\frac12\left(\frac34\right)^2=\frac9{32}}.
\]
\end{xsolution}

\paragraph{第一组参考题 15}
\begin{xsolution}
甲第一次掷出六点的概率为 $1/6$。若前两次都没有六点，则重新回到“轮到甲”的同一状态，其概率为 $(5/6)^2$。因此甲首先掷出六点的概率为
\[
  \sum_{k=0}^{\infty}
  \left(\frac56\right)^{2k}\frac16
  =\frac{1/6}{1-25/36}
  =\boxed{\frac6{11}}.
\]
\end{xsolution}

\paragraph{第一组参考题 16}
\begin{xsolution}
设年利率
\[
  r=\frac a{100}>0.
\]
存款当日发放 $1$ 万元，第二年发放 $2$ 万元，依此类推，因此以存款当日为时点，所需本金的现值为
\[
  10000\sum_{n=1}^{\infty}\frac{n}{(1+r)^{n-1}}.
\]
令 $q=(1+r)^{-1}$，则
\[
  \sum_{n=1}^{\infty}nq^{n-1}=\frac1{(1-q)^2}.
\]
故最少本金为
\[
\begin{aligned}
  P
  &=10000\frac1{(1-q)^2}\\
  &=10000\frac{(1+r)^2}{r^2}
  =\boxed{10000\frac{(100+a)^2}{a^2}\ \text{元}}.
\end{aligned}
\]
这个本金恰好使每次发放后剩余资金继续按年利率 $a\%$ 增值并支付后续全部奖学金；更少的本金其现值不足，故不能永久维持该发放规律。
\end{xsolution}

\subsection*{14.5.2 第二组参考题}

\paragraph{第二组参考题 1}
\begin{xsolution}
设函数列一致有界，即存在 $M>0$，使
\[
  \abs{f_n(x)}\leqslant M,
  \qquad x\in[a,b],\quad n\geqslant1.
\]
要证明积分数列收敛，只需证明它满足 Cauchy 准则。

给定 $\varepsilon>0$。由于 $f_n(x)$ 对每个 $x$ 收敛，定义下降数集
\[
  A_N=\{x\in[a,b]\mid
  \exists n,m\geqslant N,
  \ \abs{f_n(x)-f_m(x)}\geqslant\varepsilon\}.
\]
则 $A_1\supset A_2\supset\cdots$，且
\[
  \bigcap_{N=1}^{\infty}A_N=\varnothing.
\]
按命题 14.2.5（Lewin 引理），若
\[
  \alpha_N=\sup\{m(E)\mid E\subset A_N,\ E\text{ 为初等集}\},
\]
则 $\alpha_N\to0$。因此可取 $N$，使 $\alpha_N<\varepsilon$。

固定 $n,m\geqslant N$。函数 $\abs{f_n-f_m}$ 可积。对任意满足
\[
  0\leqslant s(x)\leqslant\abs{f_n(x)-f_m(x)}
\]
的阶梯函数 $s$，令
\[
  E=\{x\in[a,b]\mid s(x)\geqslant\varepsilon\}.
\]
则 $E$ 是初等集且 $E\subset A_N$，所以 $m(E)<\varepsilon$。在 $E$ 上有 $s\leqslant2M$，在补集上有 $s<\varepsilon$，于是
\[
  \int_a^b s(x)\,\dd x
  \leqslant2M\varepsilon+(b-a)\varepsilon.
\]
利用 Riemann 积分的 Darboux 下积分刻画，得到
\[
  \int_a^b\abs{f_n-f_m}\,\dd x
  \leqslant(2M+b-a)\varepsilon.
\]
因此
\[
  \abs{\int_a^bf_n-\int_a^bf_m}
  \leqslant(2M+b-a)\varepsilon.
\]
因 $\varepsilon$ 任意，积分数列为 Cauchy 数列，故必有有限极限。
\end{xsolution}

\paragraph{第二组参考题 2}
\begin{xsolution}
先证 $f\in R[a,b]$。给定 $\varepsilon>0$，由一致收敛可取 $N$，使
\[
  \sup_{[a,b]}\abs{f-f_N}
  <\frac{\varepsilon}{4(b-a)}.
\]
因 $f_N\in R[a,b]$，存在分划 $P$，使 $f_N$ 关于 $P$ 的 Darboux 上、下和之差小于 $\varepsilon/2$。在每个小区间 $I$ 上，
\[
  \operatorname{osc}_I f
  \leqslant\operatorname{osc}_I f_N
  +2\sup_{[a,b]}\abs{f-f_N}.
\]
乘以区间长度并求和，得
\[
  U(f,P)-L(f,P)<\frac{\varepsilon}{2}+
  \frac{\varepsilon}{2}=\varepsilon.
\]
所以 $f$ Riemann 可积。

再由
\[
  \abs{\int_a^bf_n(x)\,\dd x-
  \int_a^bf(x)\,\dd x}
  \leqslant(b-a)\sup_{[a,b]}\abs{f_n-f}
\]
以及右端趋于 $0$，得到
\[
  \boxed{
  \lim_{n\to\infty}\int_a^bf_n(x)\,\dd x
  =\int_a^bf(x)\,\dd x}.
\]
\end{xsolution}

\paragraph{第二组参考题 3}
\begin{xsolution}
不妨设所有 $f_n$ 都在 $[a,b]$ 上单调增加；单调减少时把不等号反向即可。逐点极限 $f$ 也单调增加，并且题设 $f\in C[a,b]$，故 $f$ 一致连续。

给定 $\varepsilon>0$，取分划
\[
  a=x_0<x_1<\cdots<x_m=b
\]
足够细，使
\[
  f(x_i)-f(x_{i-1})<\frac\varepsilon3,
  \qquad i=1,\ldots,m.
\]
因分点只有有限个且 $f_n(x_i)\to f(x_i)$，存在 $N$，使 $n\geqslant N$ 时对所有 $i$ 都有
\[
  \abs{f_n(x_i)-f(x_i)}<\frac\varepsilon3.
\]
任取 $x\in[x_{i-1},x_i]$。由单调性，
\[
  f_n(x_{i-1})\leqslant f_n(x)\leqslant f_n(x_i),
\]
且
\[
  f(x_{i-1})\leqslant f(x)\leqslant f(x_i).
\]
于是
\[
  f_n(x)-f(x)
  \leqslant f_n(x_i)-f(x_{i-1})
  <\frac{2\varepsilon}{3},
\]
同理
\[
  f(x)-f_n(x)<\frac{2\varepsilon}{3}.
\]
因此 $\sup_{[a,b]}\abs{f_n-f}\to0$，即一致收敛。
\end{xsolution}

\paragraph{第二组参考题 4}
\begin{xsolution}
设对每个固定 $x$，数列 $f_n(x)$ 单调增加并收敛于 $f(x)$。令
\[
  m_n=\min_{x\in[a,b]}f_n(x).
\]
因为 $f_n\leqslant f_{n+1}$，所以 $m_n$ 单调增加。又对任意固定 $x_0\in[a,b]$，
\[
  m_n\leqslant f_n(x_0)\leqslant f(x_0),
\]
故 $\{m_n\}$ 有上界，从而 $m_n\to m$。

对每个 $n$ 取 $x_n\in[a,b]$ 使 $f_n(x_n)=m_n$。由紧致性可取子列 $x_{n_k}\to x_*\in[a,b]$。固定任意 $j$，当 $n_k\geqslant j$ 时
\[
  f_j(x_{n_k})\leqslant f_{n_k}(x_{n_k})=m_{n_k}.
\]
令 $k\to\infty$，利用 $f_j$ 连续，得到
\[
  f_j(x_*)\leqslant m.
\]
再令 $j\to\infty$，得 $f(x_*)\leqslant m$。另一方面对每个 $n$，
\[
  m_n\leqslant f_n(x_*)\leqslant f(x_*),
\]
令 $n\to\infty$ 得 $m\leqslant f(x_*)$。故
\[
  f(x_*)=m,
\]
所以 $f$ 在 $[a,b]$ 上必有最小值。

若把“最小值”换成“最大值”，结论不成立。例如在 $[0,1]$ 上令
\[
  f_n(x)=\min\{x,n(1-x)\}.
\]
每个 $f_n$ 连续，并且对每个 $x$ 随 $n$ 单调增加。极限函数为
\[
  f(x)=
  \begin{cases}
    x,&0\leqslant x<1,\\
    0,&x=1.
  \end{cases}
\]
它的上确界为 $1$，但不取得最大值。

若把 $[a,b]$ 换成开区间 $(a,b)$，即使最小值结论也不成立。例如在 $(0,1)$ 上取
\[
  f_n(x)=x-\frac1n.
\]
对每个固定 $x$，数列 $\{f_n(x)\}$ 严格单调增加并收敛于 $f(x)=x$，而 $f$ 在 $(0,1)$ 上没有最小值。
\end{xsolution}

\paragraph{第二组参考题 5}
\begin{xsolution}
设
\[
  f_n'\to g
\]
在 $[a,b]$ 上一致。由一致收敛，导函数列一致有界：存在 $M>0$，使
\[
  \abs{f_n'(x)}\leqslant M
\]
对所有 $n$ 和 $x\in[a,b]$ 成立。

由 $x_n\in[a,b]$ 的紧致性，存在子列 $x_{n_k}\to x_0\in[a,b]$。

\textbf{(1)} 对任意 $x\in[a,b]$，由中值定理和 $f_{n_k}(x_{n_k})=x_{n_k}$，
\[
  \abs{f_{n_k}(x)-x_{n_k}}
  \leqslant M\abs{x-x_{n_k}}.
\]
特别
\[
  \abs{f_{n_k}(x_0)-x_0}
  \leqslant(M+1)\abs{x_{n_k}-x_0}\to0.
\]
因此锚点值 $f_{n_k}(x_0)$ 收敛。

又因 $f_n'$ 一致收敛，它一致满足 Cauchy 条件。对 $k,\ell$，
\[
\begin{aligned}
  \abs{f_{n_k}(x)-f_{n_\ell}(x)}
  &\leqslant\abs{f_{n_k}(x_0)-f_{n_\ell}(x_0)}\\
  &\quad +(b-a)\sup_{t\in[a,b]}
    \abs{f_{n_k}'(t)-f_{n_\ell}'(t)}.
\end{aligned}
\]
右端与 $x$ 无关并趋于 $0$，所以 $\{f_{n_k}\}$ 在 $[a,b]$ 上一致 Cauchy，从而一致收敛于某函数 $f$。

\textbf{(2)} 由
\[
  \abs{f_{n_k}(x_0)-x_0}\to0
\]
和一致收敛，直接得到
\[
  \boxed{f(x_0)=x_0}.
\]

\textbf{(3)} 函数 $g$ 是连续函数列 $f_n'$ 的一致极限，所以 $g\in C[a,b]$。对任意 $x\in[a,b]$，
\[
  f_{n_k}(x)-f_{n_k}(x_0)
  =\int_{x_0}^{x}f_{n_k}'(t)\,\dd t.
\]
令 $k\to\infty$，利用两边的一致收敛，得到
\[
  f(x)-f(x_0)=\int_{x_0}^{x}g(t)\,\dd t.
\]
故
\[
  f\in C^1[a,b],
  \qquad
  f'(x)=g(x)=\lim_{n\to\infty}f_n'(x).
\]
\end{xsolution}

\paragraph{第二组参考题 6}
\begin{xsolution}
由
\[
  \abs{u_n(x)}\leqslant v_n(x)
\]
可知 $v_n(x)\geqslant0$。设
\[
  V_N(x)=\sum_{n=1}^{N}v_n(x),
  \qquad
  V(x)=\sum_{n=1}^{\infty}v_n(x).
\]
每个 $V_N$ 连续，并且 $V_N(x)$ 对 $N$ 单调增加，逐点收敛于连续函数 $V$。由 Dini 定理，$V_N\to V$ 在 $[a,b]$ 上一致。因此
\[
  \sup_{x\in[a,b]}
  \sum_{n=N+1}^{\infty}v_n(x)\to0.
\]
于是
\[
  \sup_{x\in[a,b]}
  \sum_{n=N+1}^{\infty}\abs{u_n(x)}
  \leqslant
  \sup_{x\in[a,b]}
  \sum_{n=N+1}^{\infty}v_n(x)\to0.
\]
所以 $\sum u_n(x)$ 在 $[a,b]$ 上绝对一致收敛。各 $u_n$ 连续，故其和函数也连续。
\end{xsolution}

\paragraph{第二组参考题 7}
\begin{xsolution}
由于和式关于 $x$ 是 $2\pi$ 周期函数，并且取绝对值后关于 $x\mapsto2\pi-x$ 对称，只需考察 $0\leqslant x\leqslant\pi$。$x=0$ 时和为 $0$。设 $0<x\leqslant\pi$。

若 $nx\leqslant\sqrt\pi$，则
\[
  \abs{\sum_{k=1}^{n}\frac{\sin kx}{k}}
  \leqslant\sum_{k=1}^{n}\frac{kx}{k}
  =nx\leqslant\sqrt\pi.
\]

下面设 $nx>\sqrt\pi$，令
\[
  m=\left\lfloor\frac{\sqrt\pi}{x}\right\rfloor.
\]
若 $m\geqslant1$，前 $m$ 项满足
\[
  \abs{\sum_{k=1}^{m}\frac{\sin kx}{k}}
  \leqslant mx\leqslant\sqrt\pi.
\]
对尾和使用 Abel 变换。任意 $q\geqslant p$ 时，
\[
  \abs{\sum_{k=p}^{q}\sin kx}
  \leqslant\frac1{\sin(x/2)}
  \leqslant\frac\pi x,
\]
其中最后一步用了 Jordan 不等式。由于 $1/k$ 单调下降，故
\[
  \abs{\sum_{k=m+1}^{n}\frac{\sin kx}{k}}
  \leqslant\frac{\pi}{x(m+1)}.
\]
而 $m+1>\sqrt\pi/x$，所以右端小于 $\sqrt\pi$。因此
\[
  \abs{\sum_{k=1}^{n}\frac{\sin kx}{k}}<2\sqrt\pi.
\]

若 $m=0$，则 $x>\sqrt\pi$。直接对全和应用同一 Abel 估计，
\[
  \abs{\sum_{k=1}^{n}\frac{\sin kx}{k}}
  \leqslant\frac\pi x<\sqrt\pi<2\sqrt\pi.
\]
故对任意 $n$ 和实数 $x$ 均有
\[
  \boxed{
  \abs{\sum_{k=1}^{n}\frac{\sin kx}{k}}\leqslant2\sqrt\pi}.
\]
\end{xsolution}

\paragraph{第二组参考题 8}
\begin{xsolution}
对 $0<x<1$，由
\[
  \frac12\ln\frac{1+x}{1-x}
  =x+\frac{x^3}{3}+\frac{x^5}{5}+\cdots
\]
得到
\[
  \frac12\ln\frac{1+x}{1-x}-x
  =\sum_{k=1}^{\infty}\frac{x^{2k+1}}{2k+1}>0.
\]
又因 $2k+1\geqslant3$，
\[
  \sum_{k=1}^{\infty}\frac{x^{2k+1}}{2k+1}
  <\frac13\sum_{k=1}^{\infty}x^{2k+1}
  =\frac{x^3}{3(1-x^2)}.
\]
故
\[
  \boxed{
  0<\frac12\ln\frac{1+x}{1-x}-x
  <\frac{x^3}{3(1-x^2)}},
  \qquad0<x<1.
\]
\end{xsolution}

\paragraph{第二组参考题 9}
\begin{xsolution}
当 $x=0$ 时三个量都等于 $0$，所以末个不等式在 $x=0$ 只能取等号。下面证明：第一不等式对所有实数 $x$ 成立；第二个上界对所有实数 $x$ 成立，并且当 $x\ne0$ 时为严格不等式。令 $t=\abs{x}$。由幂级数展开，
\[
  \ee^x-\left(1+\frac xn\right)^n
  =\sum_{k=2}^{\infty}c_{n,k}x^k,
\]
其中当 $2\leqslant k\leqslant n$ 时
\[
  c_{n,k}
  =\frac1{k!}\left[1-\prod_{j=0}^{k-1}\left(1-\frac jn\right)\right]\geqslant0,
\]
而当 $k>n$ 时 $c_{n,k}=1/k!>0$。因此
\[
  \abs{\ee^x-\left(1+\frac xn\right)^n}
  \leqslant\sum_{k=2}^{\infty}c_{n,k}t^k
  =\ee^t-\left(1+\frac tn\right)^n.
\]
这证明了第一不等式。

对 $2\leqslant k\leqslant n$，利用 $0\leqslant j/n\leqslant1$ 以及
\[
  1-\prod_{j=0}^{k-1}(1-u_j)
  \leqslant\sum_{j=0}^{k-1}u_j,
\]
得到
\[
  c_{n,k}
  \leqslant\frac1{k!}\frac{k(k-1)}{2n}
  =\frac1{2n(k-2)!}.
\]
当 $k>n$ 时，也有
\[
  c_{n,k}=\frac1{k!}
  \leqslant\frac1{2n(k-2)!},
\]
因为 $k(k-1)\geqslant2n$。所以
\[
\begin{aligned}
  \ee^t-\left(1+\frac tn\right)^n
  &=\sum_{k=2}^{\infty}c_{n,k}t^k\\
  &\leqslant\frac1{2n}\sum_{k=2}^{\infty}
    \frac{t^k}{(k-2)!}
  =\frac{t^2\ee^t}{2n}.
\end{aligned}
\]
当 $t>0$ 时，由于至少有一个系数估计为严格不等式，故末式严格；当 $t=0$ 时取等号。因此正确结论为
\[
  \boxed{
  \abs{\ee^x-\left(1+\frac xn\right)^n}
  \leqslant \ee^{\abs{x}}-
  \left(1+\frac{\abs{x}}n\right)^n
  \leqslant\frac{x^2\ee^{\abs{x}}}{2n}},
\]
且最后一个“$\leqslant$”在 $x\ne0$ 时可以写成“$<$”。
\end{xsolution}

\paragraph{第二组参考题 10}
\begin{xsolution}
记
\[
  S_n=a_0+a_1+\cdots+a_n,
  \qquad
  \sigma_n=\frac{S_0+S_1+\cdots+S_{n-1}}n,
  \qquad n\geqslant1,
\]
且 $\sigma_n\to S$。

\textbf{(1)} 由
\[
  S_n=(n+1)\sigma_{n+1}-n\sigma_n
\]
以及 $\{\sigma_n\}$ 有界，可知 $S_n=\bigO(n)$，进而
\[
  a_n=S_n-S_{n-1}=\bigO(n).
\]
所以对 $\abs{x}<1$，级数
\[
  \sum_{n=0}^{\infty}a_nx^n
\]
绝对收敛。

\textbf{(2)} 在 $\abs{x}<1$，由绝对收敛可作级数运算。先有
\[
\begin{aligned}
  \sum_{n=0}^{\infty}a_nx^n
  &=(1-x)\sum_{n=0}^{\infty}S_nx^n.
\end{aligned}
\]
另一方面令
\[
  T_n=S_0+S_1+\cdots+S_n=(n+1)\sigma_{n+1},
\]
则
\[
  \sum_{n=0}^{\infty}S_nx^n
  =(1-x)\sum_{n=0}^{\infty}T_nx^n
  =(1-x)\sum_{n=0}^{\infty}(n+1)\sigma_{n+1}x^n.
\]
因此
\[
  \boxed{
  \sum_{n=0}^{\infty}a_nx^n
  =(1-x)^2\sum_{n=0}^{\infty}(n+1)\sigma_{n+1}x^n},
  \qquad \abs{x}<1.
\]

\textbf{(3)} 写
\[
  \sigma_{n+1}=S+\varepsilon_n,
  \qquad \varepsilon_n\to0.
\]
由 (2)，
\[
\begin{aligned}
  S(x)
  &=(1-x)^2\sum_{n=0}^{\infty}(n+1)(S+\varepsilon_n)x^n\\
  &=S+(1-x)^2\sum_{n=0}^{\infty}(n+1)\varepsilon_nx^n,
\end{aligned}
\]
其中用了
\[
  (1-x)^2\sum_{n=0}^{\infty}(n+1)x^n=1.
\]
给定 $\varepsilon>0$，取 $N$ 使 $n\geqslant N$ 时 $\abs{\varepsilon_n}<\varepsilon$。则尾项满足
\[
  (1-x)^2\sum_{n=N}^{\infty}(n+1)\abs{\varepsilon_n}x^n
  \leqslant\varepsilon.
\]
有限的前 $N$ 项乘以 $(1-x)^2$ 后在 $x\to1-$ 时趋于 $0$。故
\[
  \boxed{\lim_{x\to1-}S(x)=S}.
\]
\end{xsolution}

\paragraph{第二组参考题 11}
\begin{xsolution}
记
\[
  A(x)=\sum_{n=0}^{\infty}a_nx^n,
  \qquad
  B(x)=\sum_{n=0}^{\infty}b_nx^n.
\]
因 $a_n>0$ 且 $\sum a_n$ 在 $x=1$ 发散，所以
\[
  A(x)\to+\infty,
  \qquad x\to1-.
\]

给定 $\varepsilon>0$，由 $b_n/a_n\to A$，可取 $N$，使 $n\geqslant N$ 时
\[
  \abs{b_n-Aa_n}\leqslant\varepsilon a_n.
\]
于是
\[
\begin{aligned}
  \abs{B(x)-AA(x)}
  &\leqslant\sum_{n=0}^{N-1}\abs{b_n-Aa_n}x^n
  +\varepsilon\sum_{n=N}^{\infty}a_nx^n.
\end{aligned}
\]
除以 $A(x)>0$，得到
\[
  \abs{\frac{B(x)}{A(x)}-A}
  \leqslant\frac{C_N}{A(x)}+\varepsilon.
\]
令 $x\to1-$，第一项趋于 $0$，故上极限不超过 $\varepsilon$。因 $\varepsilon$ 任意，
\[
  \boxed{
  \lim_{x\to1-}\frac{B(x)}{A(x)}=A}.
\]
\end{xsolution}

\paragraph{第二组参考题 12}
\begin{xsolution}
设 $c_n\geqslant0$ 且
\[
  \sum_{n=0}^{\infty}c_n=A<\infty.
\]
对 $0\leqslant x<1$，
\[
  0\leqslant\sum_{n=0}^{\infty}c_nx^n\leqslant A.
\]
给定 $\varepsilon>0$，取 $N$ 使
\[
  \sum_{n>N}c_n<\varepsilon.
\]
则
\[
\begin{aligned}
  0\leqslant
  A-\sum_{n=0}^{\infty}c_nx^n
  &=\sum_{n=0}^{N}c_n(1-x^n)
   +\sum_{n>N}c_n(1-x^n)\\
  &\leqslant\sum_{n=0}^{N}c_n(1-x^n)+\varepsilon.
\end{aligned}
\]
令 $x\to1-$，有限和趋于 $0$，所以极限差不超过 $\varepsilon$。故
\[
  \boxed{
  \lim_{x\to1-}\sum_{n=0}^{\infty}c_nx^n=A}.
\]
\end{xsolution}

\paragraph{第二组参考题 13}
\begin{xsolution}
固定 $x\in[0,r)$。当 $x=0$ 时结论显然，以下设 $0<x<r$，并取 $y$ 使
\[
  x<y<r.
\]
由 Taylor 公式的积分余项，
\[
  f(x)=\sum_{k=0}^{n}\frac{f^{(k)}(0)}{k!}x^k+R_n(x),
  \qquad
  R_n(x)=\frac1{n!}\int_0^x f^{(n+1)}(t)(x-t)^n\,\dd t.
\]
因所有导函数非负，$R_n(x)\geqslant0$。作代换 $t=ux$，得到
\[
  \frac{R_n(x)}{x^{n+1}}
  =\frac1{n!}\int_0^1 f^{(n+1)}(ux)(1-u)^n\,\dd u.
\]
又因 $f^{(n+2)}\geqslant0$，函数 $f^{(n+1)}$ 单调增加；而 $ux<uy$，故
\[
  \frac{R_n(x)}{x^{n+1}}
  \leqslant
  \frac1{n!}\int_0^1 f^{(n+1)}(uy)(1-u)^n\,\dd u
  =\frac{R_n(y)}{y^{n+1}}.
\]
同时 Taylor 公式中的各项均非负，所以
\[
  0\leqslant R_n(y)\leqslant f(y).
\]
因此
\[
  0\leqslant R_n(x)
  \leqslant f(y)\left(\frac xy\right)^{n+1}\to0.
\]
于是对每个 $x\in[0,r)$，
\[
  \boxed{
  f(x)=\sum_{n=0}^{\infty}\frac{f^{(n)}(0)}{n!}x^n}.
\]
\end{xsolution}

\paragraph{第二组参考题 14}
\begin{xsolution}
定义生成函数
\[
  f(x)=\sum_{n=0}^{\infty}a_nx^n,
  \qquad a_0=a_1=1,
  \qquad a_n=a_{n-1}+a_{n-2}\ (n\geqslant2).
\]
在收敛半径内，
\[
\begin{aligned}
  f(x)-1-x
  &=\sum_{n=2}^{\infty}a_nx^n\\
  &=x\sum_{n=2}^{\infty}a_{n-1}x^{n-1}
   +x^2\sum_{n=2}^{\infty}a_{n-2}x^{n-2}\\
  &=x(f(x)-1)+x^2f(x).
\end{aligned}
\]
故
\[
  \boxed{f(x)=\frac1{1-x-x^2}}.
\]

令
\[
  \varphi=\frac{1+\sqrt5}{2},
  \qquad
  \psi=\frac{1-\sqrt5}{2}.
\]
因为
\[
  1-x-x^2=(1-\varphi x)(1-\psi x),
\]
作部分分式展开得
\[
  \frac1{1-x-x^2}
  =\frac1{\sqrt5}
  \left(\frac{\varphi}{1-\varphi x}
       -\frac{\psi}{1-\psi x}\right).
\]
比较 $x^n$ 的系数，
\[
  \boxed{
  a_n=\frac{\varphi^{n+1}-\psi^{n+1}}{\sqrt5}}.
\]
由于 $\abs{\psi}<1<\varphi$，
\[
  \frac{a_{n+1}}{a_n}
  =\varphi\,
  \frac{1-(\psi/\varphi)^{n+2}}
       {1-(\psi/\varphi)^{n+1}}
  \to\boxed{\varphi=\frac{1+\sqrt5}{2}}.
\]
这就是该数列的增长数。
\end{xsolution}

\paragraph{第二组参考题 15}
\begin{xsolution}
利用广义二项式展开，在 $\abs{x}<1$ 内有
\[
  (1+x)^\alpha
  =\sum_{k=0}^{\infty}{\alpha\choose k}x^k,
  \qquad
  (1+x)^\beta
  =\sum_{j=0}^{\infty}{\beta\choose j}x^j.
\]
两级数绝对收敛，可以作 Cauchy 乘积：
\[
  (1+x)^\alpha(1+x)^\beta
  =\sum_{n=0}^{\infty}
  \left(\sum_{k=0}^{n}{\alpha\choose k}{\beta\choose n-k}\right)x^n.
\]
另一方面
\[
  (1+x)^\alpha(1+x)^\beta=(1+x)^{\alpha+\beta}
  =\sum_{n=0}^{\infty}{\alpha+\beta\choose n}x^n.
\]
由幂级数展开的惟一性，比较 $x^n$ 系数得到 Vandermonde 恒等式
\[
  \boxed{
  \sum_{k=0}^{n}{\alpha\choose k}{\beta\choose n-k}
  ={\alpha+\beta\choose n}}.
\]

令 $\alpha=\beta=n$，并利用
\[
  {n\choose n-k}={n\choose k},
\]
便得到
\[
  \boxed{
  \sum_{k=0}^{n}{n\choose k}^2={2n\choose n}}.
\]
\end{xsolution}
